% concatenated from Cluster_Size_dRF_arXiv.tex
\documentclass[aap,preprint]{imsart}

%% Packages
\RequirePackage{amsthm,amsmath,amsfonts,amssymb}
\RequirePackage[numbers,sort&compress]{natbib}
%\RequirePackage[authoryear]{natbib}%% uncomment this for author-year citations
\RequirePackage[colorlinks,citecolor=blue,urlcolor=blue]{hyperref}
\RequirePackage{graphicx}
\usepackage{euscript,bbm,xcolor,epstopdf}
\usepackage{mathrsfs}
%
%\RequirePackage[colorlinks,citecolor=blue,urlcolor=blue]{hyperref}
%\usepackage[colorlinks,citecolor=blue,urlcolor=blue]{hyperref}

\usepackage{bm,euscript,bbm}
\usepackage{verbatim}
\usepackage{multirow}
\usepackage{epstopdf}
\usepackage{color}
\usepackage{pstricks,pst-node,pst-plot}
\usepackage{rotating}
\usepackage{enumitem}
\usepackage{booktabs}

\startlocaldefs
%%%%%%%%%%%%%%%%%%%%%%%%%%%%%%%%%%%%%%%%%%%%%%
%%                                          %%
%% Uncomment next line to change            %%
%% the type of equation numbering           %%
%%                                          %%
%%%%%%%%%%%%%%%%%%%%%%%%%%%%%%%%%%%%%%%%%%%%%%
%\numberwithin{equation}{section}
%%%%%%%%%%%%%%%%%%%%%%%%%%%%%%%%%%%%%%%%%%%%%%
%%                                          %%
%% For Axiom, Claim, Corollary, Hypothesis, %%
%% Lemma, Theorem, Proposition              %%
%% use \theoremstyle{plain}                 %%
%%                                          %%
%%%%%%%%%%%%%%%%%%%%%%%%%%%%%%%%%%%%%%%%%%%%%%
\theoremstyle{plain}
\newtheorem{axiom}{Axiom}
\newtheorem{claim}[axiom]{Claim}
\newtheorem{theorem}{Theorem}[section]
\newtheorem{lemma}[theorem]{Lemma}
\newtheorem{corollary}[theorem]{Corollary}
\newtheorem{proposition}[theorem]{Proposition}

%%%%%%%%%%%%%%%%%%%%%%%%%%%%%%%%%%%%%%%%%%%%%%
%%                                          %%
%% For Assumption, Definition, Example,     %%
%% Notation, Property, Remark, Fact         %%
%% use \theoremstyle{definition}            %%
%%                                          %%
%%%%%%%%%%%%%%%%%%%%%%%%%%%%%%%%%%%%%%%%%%%%%%
\theoremstyle{definition}
\newtheorem{definition}[theorem]{Definition}
\newtheorem{example}[theorem]{Example}
\newtheorem{remark}[theorem]{Remark}
\newtheorem{assumption}[theorem]{Assumption}
\newtheorem*{fact}{Fact}
%%%%%%%%%%%%%%%%%%%%%%%%%%%%%%%%%%%%%%%%%%%%%%
%% Please put your definitions here:        %%
%%%%%%%%%%%%%%%%%%%%%%%%%%%%%%%%%%%%%%%%%%%%%%

\def\R{{\mathbb R}}
\def\E{{\mathbb E}}
\def\P{{\mathbb P}}
\def\Z{{\mathbb Z}}

\endlocaldefs

\begin{document}
	
	\begin{frontmatter}
		\title{Cluster size distributions of discrete random fields}
		%\title{A sample article title with some additional note\thanksref{t1}}
		\runtitle{Cluster size distribution}
		%\thankstext{T1}{A sample additional note to the title.}
		
		\begin{aug}
			%%%%%%%%%%%%%%%%%%%%%%%%%%%%%%%%%%%%%%%%%%%%%%%
			%% Only one address is permitted per author. %%
			%% Only division, organization and e-mail is %%
			%% included in the address.                  %%
			%% Additional information such as            %%
			%% identifying the corresponding author must %%
			%% be included in in the Acknowledgments     %%
			%% section if necessary.                     %%
			%% ORCID can be inserted by command:         %%
			%% \orcid{0000-0000-0000-0000}               %%
			%%%%%%%%%%%%%%%%%%%%%%%%%%%%%%%%%%%%%%%%%%%%%%%
			\author{\fnms{Dan}~\snm{Cheng} \ead[label=e1]{chengdan@asu.edu}}
			\and
			\author{\fnms{John}~\snm{Ginos}\ead[label=e2]{jginos@asu.edu}}
			%%%%%%%%%%%%%%%%%%%%%%%%%%%%%%%%%%%%%%%%%%%%%%
			%% Addresses                                %%
			%%%%%%%%%%%%%%%%%%%%%%%%%%%%%%%%%%%%%%%%%%%%%%
			\address{School of Mathematical and Statistical Sciences, Arizona State University\\ \printead{e1,e2}}
			
		\end{aug}
		
		\begin{abstract}
			We study discrete random fields $X_t$ indexed by the integer lattice $\Z^d$. For a fixed threshold $u$, the excursion set $\{t \in \Z^d : X_t > u\}$ decomposes into connected components, called excursion clusters, whose sizes are the numbers of lattice sites they contain. Our objective is to characterize the probability distribution of cluster sizes.
			
			For stationary random fields, we define the size distribution of a typical cluster by assigning each cluster a unique root and derive an exact representation in terms of rooted cluster intensities. To address nonstationary fields, we introduce a peak-based cluster size distribution, which characterizes the distribution of cluster sizes conditional on the presence of a local maximum above $u$. It is shown that, under suitable conditions, the difference between the exact and peak-based distributions decreases exponentially as $u\to\infty$. We further develop guarded spatial Monte Carlo estimators when direct evaluation of intensities of large clusters is infeasible. Finally, we derive exponential bounds for large clusters at a fixed threshold, as well as Gaussian large deviation asymptotics at a fixed cluster size as $u\to\infty$. The framework applies to both Gaussian and non-Gaussian fields and provides a probabilistic foundation for quantifying spatial extent in discretely sampled data.
		\end{abstract}
		
		\begin{keyword}[class=MSC]
			\kwd{60G60}
			\kwd{60G15}
			\kwd{60G70} 
		\end{keyword}
		
		\begin{keyword}
			\kwd{Discrete random fields}
			\kwd{excursion sets}
			\kwd{cluster size}
			\kwd{connected component}
			\kwd{local maxima}
			\kwd{spatial Monte Carlo}
			\kwd{Gaussian large deviations}
		\end{keyword}
		
	\end{frontmatter}
	%%%%%%%%%%%%%%%%%%%%%%%%%%%%%%%%%%%%%%%%%%%%%%
	%% Please use \tableofcontents for articles %%
	%% with 50 pages and more                   %%
	%%%%%%%%%%%%%%%%%%%%%%%%%%%%%%%%%%%%%%%%%%%%%%
	%\tableofcontents
	
	\section{Introduction}	
	Random fields play a central role in modern spatial statistics, stochastic modeling, and signal detection. They provide a natural framework for characterizing spatially correlated data across a wide range of scientific disciplines, including medical imaging, geoscience, environmental studies, and cosmology \cite{AT07,Cheng:2020,Chiles:2012,Goeman:2023,MP11,Taylor:2007}.
	When a random field is observed over a discrete lattice such as $\Z^d$, the resulting process is referred to as a \textit{discrete random field}. A fundamental problem in the analysis of such fields concerns the geometric and topological properties of regions where the field exceeds a fixed threshold level 
	$u$. These regions, commonly called excursion sets, decompose into random connected components, or clusters, formed by neighboring lattice sites above the threshold.
	
	Understanding the statistical distribution of cluster sizes in random fields is fundamental for both theory and applications. In spatial statistics and random field theory, the cluster size distribution under a null (pure noise) model provides a benchmark for assessing the significance of observed spatial patterns. For example, in neuroimaging, random field theory underlies cluster-based inference procedures for detecting activation regions in fMRI data \cite{Bansal:2018,Chumbley:2010,Friston:1994,Goeman:2023,Poline:1997,Worsley:1996a,Zhang:2009}. Similar principles arise in geoscience, genomics and cosmology, where cluster size distributions describe the spatial extent of anomalies in correlated data  \cite{Adler81,Bardeen:1985,Benjamini:2019,Lindgren82,Vanmercke:2010}. Recent work in climate and environmental statistics has emphasized the importance of spatial extent inference, which aims to quantify the areal footprint of extreme events \cite{Tan:2021,Zhong:2025}. In parallel, inference for threshold exceedances in random fields has been applied to region detection \cite{Benjamini:2019} and spatiotemporal extreme value analysis \cite{Fondeville:2018,Huser:2014}. 
	
	While the geometry of excursion sets for continuous Gaussian fields has been extensively studied through tools such as the Gaussian kinematic formula and Euler characteristic heuristics \cite{AT07,Taylor:2007}, exact cluster size distributions are typically intractable in continuous settings. In contrast, for discrete random fields on lattices, it is possible to formulate and compute cluster size distributions explicitly under suitable assumptions. The discrete setup allows a combinatorial characterization of connected components, leading to exact expressions in low dimensions and numerically tractable formulations in higher dimensions.
	
	In this paper, we investigate the cluster size distribution of discrete stationary random fields $\{X_t: t\in \Z^d\}$ defined on the $d$-dimensional integer lattice $\Z^d$. We focus on the probabilistic structure of clusters that arise from thresholding the field at a fixed level $u\in \R$. Denote the excursion set above $u$ by
	\[
	E_u=\{t\in\mathbb Z^d:X_t>u\}.
	\]
	For $t\in E_u$, let \(C_u(t)\) denote the connected component
	of \(E_u\) containing \(t\); set \(C_u(t)=\varnothing\) when
	\(t\notin E_u\). The size $|C_u(t)|$ is defined as the number of lattice points in $C_u(t)$. Throughout this paper, we will make use of the following finite cluster assumption.
	
	\begin{assumption}[Finite clusters]\label{ass:finite-cluster}
		At the threshold $u$ under consideration,
		\begin{equation*}
			\mathbb P\left(
			|C_u(t)|<\infty \text{ for every }t\in E_u
			\right)=1.
		\end{equation*}
	\end{assumption}

	Thus every excursion cluster considered in this paper is finite
	almost surely. The assumption excludes percolating excursion components. If it holds at $u$, then it automatically holds for $u'>u$. For a stationary field, Assumption~\ref{ass:finite-cluster}
	is equivalent to $\mathbb P\bigl(|C_u(t_0)|=\infty\bigr)=0$ for some $t_0$. In the stationary case, due to the invariant translation property, we may simply denote a cluster $C_u(t_0)$ by $C_u$.
	
	Illustrations of clusters and their sizes in one and two dimensions are provided in Figures \ref{fig:simul-1D} and \ref{fig:simul-2D}. We first derive in Section \ref{sec:exact-cluster} the \emph{exact cluster size distribution} for stationary random fields (see also Definition \ref{def:cluster-size}):
	\begin{equation*}
		\begin{split}
			\P( |C_u|=k \ | \ \text{a cluster $C_u$ above $u$ exists}), \quad k=1,2,\ldots.
		\end{split}
	\end{equation*}
	The exact distributional formulae are shown in Theorem \ref{thm:CSD_d} and Corollary \ref{cor:CSD_1D}.
	
	Motivated by recent advances in peak inference for random fields \citep{CS15,CS18,CS17} and peaks-over-threshold methods \citep{Fondeville:2018,Huser:2014}, we further introduce and develop the \textit{peak-based cluster size distribution} in Section \ref{sec:peak-based-cluster}. This characterizes the distribution of cluster sizes conditional on the presence of a local maximum above $u$, and can be formally written as the conditional probability (see also Definition \ref{def:peak-cluster-size})
	\[
	\P(|C_u(t)|=k \ | \ t \ \text{\rm is a local maximum and } X_t>u), \quad k=1,2,\ldots.
	\]
	The formula is derived in Theorem \ref{thm:peak-CSD}. It links local peak behavior with cluster extent, thereby connecting spatial extreme value theory with classical random field analysis. It is particularly valuable in nonstationary settings, where local variations in the covariance structure influence both the shape and spatial extent of clusters. In contrast, obtaining exact cluster size distributions in the nonstationary case appears analytically intractable. 
	
	At high thresholds, excursion clusters typically contain a unique local maximum, suggesting that the peak-based and exact cluster size distributions should be close for stationary fields. In Section \ref{sec:high-u}, we quantify this approximation by bounding the total variation distance between the two distributions. For stationary Gaussian fields, we further show that this distance decays exponentially in $u^2$ as $u\to\infty$.
	
	A further contribution is a Monte Carlo framework for estimating both cluster size distributions in Section \ref{sec:MC}. To estimate probabilities for cluster sizes up to a cutoff $K$, we simulate the field on an enlarged window containing a guard band around the scoring window. This guard band determines exactly whether each scored site belongs to a cluster of size $1,\ldots,K$ or to a larger cluster, thereby eliminating finite-window bias. We prove that the resulting spatial estimators of the unnormalized cluster weights are unbiased and strongly consistent, derive a computable upper bound for the normalization error caused by clusters larger than $K$, and obtain an exact variance formula that explicitly accounts for the size and geometry of the scoring window.
	
	Moreover, we analyze two distinct asymptotic regimes in Section \ref{sec:large-deviations}. First, with the threshold $u$ fixed and the cluster size $k\to \infty$, we establish exponential upper bounds for the tails of both the exact and peak-based cluster size distributions. For Gaussian fields, we provide an explicit sufficient condition based on covariance summability. Second, with $k$ fixed and $u\to\infty$, finite-dimensional Gaussian large deviation theory gives exact logarithmic decay rates, with speed $u^2$, for both distributions. 
	
	In Section~\ref{sec:simulation}, we present numerical experiments in one and two dimensions to illustrate the theoretical formulas and spatial Monte Carlo estimators, and to examine how spatial dependence and the choice of neighborhood connectivity affect the cluster size distributions.
	
	In many applications, data are observed on discrete grids rather than as continuous fields. The framework developed in this paper thus provides a unified theoretical foundation for cluster modeling and inference in discrete random fields. Beyond its intrinsic theoretical interest, the results have broad practical implications for signal detection and spatial extent inference in correlated lattice data. Applications include neuroimaging \cite{Bansal:2018,Chumbley:2010,Friston:1994,Goeman:2023,Worsley:1996a}, geophysical anomaly detection \cite{Chiles:2012,Vanmercke:2010}, and climate data analysis \cite{Cooley:2007,Davison:2012}, where discretely sampled random fields are ubiquitous.
	
	Finally, while Gaussian random fields serve as a convenient and widely used reference model, the general framework developed here does not require Gaussianity. As will become clear from the arguments below, the derivations depend only on the joint dependence structure of the field. Consequently, the methodology applies to broad classes of non-Gaussian random fields, provided their finite-dimensional distributions are well defined.
	
	
	\section{The exact cluster size distribution}\label{sec:exact-cluster}
	
	We begin directly with a stationary random field on $\Z^d$. The key device is
	to assign each finite excursion cluster a canonical root. This makes it
	possible to define a translation-invariant distribution over clusters and to
	derive the exact cluster size distribution in arbitrary dimension. We then
	justify rigorously its empirical interpretation through root counts in
	expanding observation windows, including the effect of discarding clusters
	near the observation boundary. 
	
	\subsection{General $d$-dimensional formulation}
	
	Let $\{X_t:t\in\Z^d\}$ be a real-valued stationary random field, let
	$u\in\R$ be fixed, and suppose Assumption~\ref{ass:finite-cluster} holds.  Recall
	that
	\[
	E_u=\{t\in\Z^d:X_t>u\}
	\]
	and that $C_u(t)$ denotes the connected component of $E_u$ containing $t$,
	with $C_u(t)=\varnothing$ when $t\notin E_u$.
	
	\subsubsection{Connectivity, roots, and rooted cluster shapes}
	
	Denote the origin by $o=(0,\ldots,0)$, and let
	$e_1=(1,0,\ldots,0),\ldots,e_d=(0,\ldots,0,1)$ be the standard basis vectors.
	We consider either of the following two neighborhood systems.  The
	\emph{nearest neighbors} of $t\in\Z^d$ are
	\begin{equation*}
		\{t'\in\Z^d:\|t'-t\|=1\}
		=\{t\pm e_i:i=1,\ldots,d\},
	\end{equation*}
	so each site has $2d$ nearest neighbors.  The \emph{Moore neighbors} of $t$ are
	\begin{equation*}
		\{t'\in\Z^d:\|t'-t\|_\infty=1\},
	\end{equation*}
	where $\|(t_1,\ldots,t_d)\|_\infty=\max_{1\le i\le d}|t_i|$.  Thus, the Moore
	neighborhood also includes diagonal sites and contains $3^d-1$ sites. Figure \ref{fig:neighbor} illustrates these two neighborhood systems for $d=2$. We fix
	one of these two neighborhood systems throughout each argument.
	
	\begin{figure}[t]
		\begin{center}
			\begin{tabular}{cc}
				\includegraphics[trim=0 160 0 0,clip,width=1.5in]{Neighbor_4.png} &
				\includegraphics[trim=0 160 0 0,clip,width=1.5in]{Neighbor_8.png} \\
				{ Nearest-neighbor connectivity} & { Moore connectivity}
			\end{tabular}
			\caption{\label{fig:neighbor} Two neighborhood systems in $\Z^2$.
				}
		\end{center}
	\end{figure}
	
	For a finite set $D\subset\Z^d$, let $\mathcal{N}(D)$ denote its
	\emph{exterior neighbor set}:
	\begin{equation}\label{eq:neighbor}
		\mathcal{N}(D)=
		\begin{cases}
			\{t'\in\Z^d\setminus D:\|t'-t\|=1
			\text{ for some }t\in D\}, & \text{nearest neighbors},\\
			\{t'\in\Z^d\setminus D:\|t'-t\|_\infty=1
			\text{ for some }t\in D\}, & \text{Moore neighbors}.
		\end{cases}
	\end{equation}
	For a singleton, we write $\mathcal{N}(t)=\mathcal{N}(\{t\})$.  We also use
	$X_D>u$ to mean $X_t>u$ for every $t\in D$, and define $X_D\le u$
	analogously.
	
	To select one representative location for each cluster, introduce the
	lexicographic order on $\Z^d$.  For distinct
	$t=(t_1,\ldots,t_d)$ and $t'=(t'_1,\ldots,t'_d)$, write
	\[
	t<_{\rm lex}t'
	\quad\Longleftrightarrow\quad
	\text{for some $j$, $t_j<t'_j$ and $t_i=t'_i$ for every $i<j$}.
	\]
	Thus the first coordinate where the two vectors differ determines the ordering. For example, in $\Z^2$,
	\[
	(0,0)<_{\rm lex}(0,1)<_{\rm lex}(1,-1)<_{\rm lex}(1,0)<_{\rm lex}(1,1).
	\]
	Every finite nonempty subset of $\Z^d$ has a unique lexicographically smallest
	point.  We call this point the \emph{root} of the cluster.  For every
	$t\in\Z^d$, define the root event
	\begin{equation}\label{eq:translated-root-event}
		R_u(t)
		=
		\{X_t>u,\ t<_{\rm lex}s
		\text{ for every }s\in C_u(t)\setminus\{t\}\}.
	\end{equation}
	Thus, $R_u(t)$ is the event that $t$ is the root of its excursion cluster.
	In particular, $R_u(o)$ denotes the event that the origin is a cluster root.
	
	
	For $k\ge1$, define the deterministic collection of size-$k$ connected shapes rooted at $o$ by
	\begin{equation*}
		\mathcal{C}_k^{\rm root}
		=
		\left\{
		D\subset\Z^d:
			|D|=k,\ D\text{ is connected with root at } o
		\right\}.
	\end{equation*}
	
	\begin{figure}[t]
		\begin{center}
			\begin{tabular}{c}
				\includegraphics[trim=8 8 25 5,clip,width=5in]{Toy_1D.png}\\[.5em]
				{(a) Simulated Gaussian process and clusters above $u=0.5$}
				\vspace{4mm}\\
				\includegraphics[trim=5 10 50 5,clip,width=3in]
				{empirical_CSD_1d_Gauss_u_eq0p5.png}\\[.5em]
				{(b) Empirical cluster size distribution}
			\end{tabular}
			\caption{\label{fig:simul-1D}
				Panel (a) is a scatter plot of a simulated realization of a centered
				stationary Gaussian process with covariance
				$C(t,s)=e^{-(t-s)^2}$.  Exceedances above $u=0.5$ and the corresponding
				clusters are marked by consecutive red dots along $\Z$.  This realization
				has 14 observed clusters: 7 of size one, 5 of size two, 1 of size three,
				and 1 of size four.  Panel (b) displays the empirical cluster size
				distribution obtained from repeated simulations.}
		\end{center}
	\end{figure}
	
	\noindent\textbf{Examples of rooted shapes.}
	In one dimension, the nearest and Moore neighbors coincide. For
	every $k\geq1$, there is exactly one rooted connected shape,
	\begin{equation}\label{eq:root-shapes-1d}
		D_k=\{0,1,\ldots,k-1\},
		\qquad
		\mathcal C_k^{\rm root}=\{D_k\},
		\qquad
		\mathcal N(D_k)=\{-1,k\}.
	\end{equation}
	In two dimensions with nearest-neighbor connectivity, the first three rooted
	collections are
	\begin{equation}\label{eq:root3}
		\begin{split}
			\mathcal{C}_1^{\rm root}
			&=\bigl\{\{o\}\bigr\},\\
			\mathcal{C}_2^{\rm root}
			&=\bigl\{\{o,e_1\},\{o,e_2\}\bigr\},\\
			\mathcal{C}_3^{\rm root}
			&=\bigl\{
			\{o,e_1,2e_1\},
			\{o,e_1,e_1+e_2\},
			\{o,e_1,e_1-e_2\},\\
			&\hspace{18mm}
			\{o,e_1,e_2\},
			\{o,e_2,2e_2\},
			\{o,e_2,e_1+e_2\}
			\bigr\}.
		\end{split}
	\end{equation}
	Under Moore connectivity, $\mathcal C_2^{\rm root}$ additionally contains
	$\{o,e_1+e_2\}$ and $\{o,e_1-e_2\}$, and
	$\mathcal C_3^{\rm root}$ contains further shapes because diagonal sites are
	also adjacent.
	
	For $D\in\mathcal{C}_k^{\rm root}$, the event that $D$ is exactly the
	excursion cluster rooted at $o$ is
	\begin{equation}\label{eq:C_u(o)}
	\{C_u(o)=D\}=\{X_D>u,\ X_{\mathcal{N}(D)}\le u\}.
    \end{equation}
	
	\begin{definition}\label{def:cluster-size} For a discrete stationary random field, the exact distribution of the cluster size $S_u$ is defined as 
		\begin{equation}\label{eq:CSD_def}
			\P(S_u=k)
			= \P(|C_u|=k \ | \ \text{a cluster $C_u$ above $u$ exists}), \quad k=1,2,\ldots,
		\end{equation}
		provided the conditional event has a positive probability. 
	\end{definition}
	More precisely, the conditional event refers to selecting a cluster under a sampling rule that assigns equal weight to every excursion cluster. Figures \ref{fig:simul-1D} and \ref{fig:simul-2D} provide visual illustrations for cases $d=1$ and $d=2$, respectively. By stationarity, Definition~\ref{def:cluster-size} does not depend on a
	particular lattice location. 
	
	\begin{figure}[t]
		\begin{center}
			\begin{tabular}{cc}
				\includegraphics[trim=20 20 0 0,clip,width=1.5in]{Toy_2D_4.png} &
				\includegraphics[trim=0 0 0 0,clip,width=3in]
				{SqrExpCovGRF_4N_CSD_2D_u_eq_0p5.png}\\[.5em]
				 (a) Clusters (nearest neighbors) &
				 (b) Empirical distribution (nearest neighbors)
			\end{tabular}
			
			\vspace{3.9mm}
			
			\begin{tabular}{cc}
				\includegraphics[trim=20 20 0 0,clip,width=1.5in]{Toy_2D_8.png} &
				\includegraphics[trim=0 0 0 0,clip,width=3in]
				{SqrExpCovGRF_8N_CSD_2D_u_eq_0p5.png}\\[.5em]
				 (c) Clusters (Moore neighbors) &
				 (d) Empirical distribution (Moore neighbors)
			\end{tabular}
			\caption{\label{fig:simul-2D}
				Clusters and empirical cluster size distributions for an isotropic
				Gaussian field on $\Z^2$ with covariance
				$C(t,s)=e^{-\|t-s\|^2}$ and threshold $u=0.5$.  Panels (a) and (b)
				use nearest-neighbor connectivity; panels (c) and (d) use Moore
				connectivity.  Distinct clusters above $u$ are highlighted in panels
				(a) and (c), and panels (b) and (d) show the corresponding empirical
				distributions from repeated simulations.}
		\end{center}
	\end{figure}
	
	\subsubsection{Exact distribution in $\Z^d$}
	
	Let $w_k$ denote the probability that there exists a size-$k$ cluster above $u$ with root at $o$, that is,
	\begin{equation}\label{eq:w_d}
		\begin{split}
			w_k
			&:=\sum_{D\in\mathcal{C}_k^{\rm root}}
			\P\bigl(X_D>u,\ X_{\mathcal{N}(D)}\le u\bigr),
			\qquad k=1,2,\ldots.
		\end{split}
	\end{equation}
	The quantity $w_k$ is the intensity, or expected number per lattice site, of
	size-$k$ clusters when each cluster is counted once at its root. For fixed $k$, the events
	\[
	\{X_D>u,\ X_{\mathcal{N}(D)}\le u\},
	\qquad D\in\mathcal{C}_k^{\rm root},
	\]
	are pairwise disjoint, and their union is
	$\{R_u(o),\ |C_u(o)|=k\}$.  Hence we have
	\begin{equation}\label{eq:w_d_2}
		\begin{split}
			w_k =\P\bigl(R_u(o),\ |C_u(o)|=k\bigr),
			\qquad k=1,2,\ldots.
		\end{split}
	\end{equation}
	
	\begin{theorem}\label{thm:CSD_d}
		Let $\{X_t:t\in\Z^d\}$ be stationary, let $u\in\R$ satisfy
		Assumption~\ref{ass:finite-cluster}, and suppose
		$\rho_u=\P(R_u(o))>0$.  Then
		\begin{equation}\label{eq:CSD_d}
			\P(S_u=k)=\frac{w_k}{\rho_u},
			\qquad
			\rho_u=\sum_{j\geq1}w_j,
			\qquad k\geq1,
		\end{equation}
		where $w_k$ is given by \eqref{eq:w_d}.
	\end{theorem}
	
	\begin{proof}
		By stationarity, without loss of generality, we assume that the cluster under consideration, denoted $C_u$, has root at the origin $o$. Recall that $R_u(o)$ denotes the event that a cluster above $u$ has root at $o$. Therefore,
		\begin{equation}\label{eq:CSD_root_interpretation}
			\begin{split}
				\P(S_u=k)
				&= \P\bigl(|C_u|=k \,\big|\, \text{there exists a cluster $C_u$ above $u$ with root at $o$}\bigr)\\
				&= \P\bigl(|C_u(o)|=k\mid R_u(o)\bigr).
			\end{split}
		\end{equation}
		It then follows from
		\eqref{eq:w_d_2} that
		\[
		\P(S_u=k)
		=\frac{\P(R_u(o),|C_u(o)|=k)}{\P(R_u(o))}
		=\frac{w_k}{\rho_u}.
		\]
		By Assumption~\ref{ass:finite-cluster}, on $R_u(o)$ the cluster $C_u(o)$ has some finite size.  Hence
		\[
		R_u(o)=\bigsqcup_{j\ge1}
		\{R_u(o),\ |C_u(o)|=j\},
		\]
		together with \eqref{eq:w_d_2}, we obtain
		\[
		\rho_u = \P(R_u(o)) = \sum_{j=1}^\infty \P\bigl(R_u(o),\ |C_u(o)|=j\bigr) = \sum_{j=1}^\infty w_j,
		\]
		completing the proof.
	\end{proof}
	
	\begin{remark}
		It is seen from \eqref{eq:CSD_root_interpretation} that, the exact cluster size distribution can be precisely written by the
		root convention:
		\begin{equation*}
			\P(S_u=k) = \P\bigl(|C_u(o)|=k\mid R_u(o)\bigr),
		\end{equation*}
		provided $\P(R_u(o))>0$.
	\end{remark}
	
	\subsection{Consistency of empirical cluster frequencies}
	\label{sec:empirical-cluster-counts}
	We now demonstrate that Theorem~\ref{thm:CSD_d} is consistent with the empirical cluster size distribution computed in simulations when expanding observation windows.
	
	Recall from \eqref{eq:translated-root-event} that
	$R_u(t)$ is the event that $t$ is a cluster root.  Define
	\begin{equation}\label{eq:empirical-root-indicators}
		Y_{t,k}=\mathbbm{1}\{R_u(t),|C_u(t)|=k\},
		\qquad
		Y_t=\mathbbm{1}\{R_u(t)\}.
	\end{equation}
	The indicator $Y_{t,k}$ contributes one
	precisely when $t$ is the unique root of a full-lattice excursion cluster of
	size $k$.  For a finite window $\Lambda$, the \textit{ideal root counts} are
	\begin{equation}\label{eq:ideal-root-counts}
		N_k^*(\Lambda)=\sum_{t\in\Lambda}Y_{t,k},
		\qquad
		N^*(\Lambda)=\sum_{t\in\Lambda}Y_t.
	\end{equation}
	 Since every finite cluster has exactly one root,
	$N_k^*(\Lambda)$ therefore counts the size-$k$ full-lattice clusters whose
	roots lie in $\Lambda$ (part of the cluster may lie outside $\Lambda$), with each such cluster counted exactly once. Similarly, $N^*(\Lambda)$ is the total number of full-lattice excursion clusters, of all sizes, whose roots lie in $\Lambda$. Under Assumption~\ref{ass:finite-cluster}, one has $N^*(\Lambda)=\sum_{k=1}^{\infty}N_k^*(\Lambda)$.
	
	
	We next define the clusters that can be identified using only the observed
	values $\{X_t:t\in\Lambda\}$. First, let $\mathscr{C}_u(\Lambda)$ denote the collection of observed clusters restricted on $\Lambda$:
	\begin{equation}\label{eq:clusters-restriction}
		\mathscr{C}_u(\Lambda)
		=
		\{\text{\rm connected components of } E_u\cap\Lambda\}.
	\end{equation}
	Thus, every $D\in \mathscr{C}_u(\Lambda)$ satisfies
	$D\subset\Lambda$ and is connected through lattice sites in $\Lambda$, although it may be the observed portion of a cluster
	that continues outside $\Lambda$.
	
	Second, call an observed component $D$ an \emph{interior cluster} (with
	respect to $\Lambda$) if its entire exterior neighbor set is contained in the observation window $\Lambda$. The collection of interior clusters is
	\begin{equation}\label{eq:observable-interior-components}
		\mathscr{C}_u^\circ(\Lambda)
		=
		\{D\in \mathscr{C}_u(\Lambda):\mathcal N(D)\subset\Lambda\}.
	\end{equation}
	This interior condition makes $D$ a genuine full-lattice excursion cluster.
	Indeed, $\mathcal N(D)\subset\Lambda$ implies that every
	exterior neighbor of $D$ is observed and is at or below $u$.  Hence the
	cluster cannot continue outside $D$. Define the \textit{observable counts} or \textit{empirical counts}
	\begin{equation}\label{eq:observable-cluster-counts}
	\begin{split}
		N_k^\circ(\Lambda)
		&=\sum_{D\in \mathscr{C}_u^\circ(\Lambda)}
		\mathbbm{1}\{|D|=k\}=\#\{\text{size-$k$ interior clusters in $\Lambda$}\},\\
		\qquad
		N^\circ(\Lambda)&=| \mathscr{C}_u^\circ(\Lambda)|=\#\{\text{all interior clusters in $\Lambda$}\}.
		\end{split}
	\end{equation}
	
	Let $d_G$ be the graph distance associated with the fixed neighborhood
	system. For an integer $r\geq0$, write
	\begin{equation}\label{eq:empirical-window-boundary}
		\partial_r\Lambda
		=\{t\in\Lambda:d_G(t,\Lambda^c)\leq r\}.
	\end{equation}
	We use nested rectangular windows centered at the origin of the form
	\begin{equation}\label{eq:empirical-rectangular-windows}
		\Lambda_n
		=\prod_{i=1}^d\{-m_{i,n},\ldots,m_{i,n}\},
		\qquad m_{i,n}\uparrow\infty\quad(i=1,\ldots,d).
	\end{equation}
	In particular,
	for each fixed $r$,
	\begin{equation}\label{eq:empirical-boundary-negligible}
		\frac{|\partial_r\Lambda_n|}{|\Lambda_n|}\longrightarrow0.
	\end{equation}
	
	\begin{theorem}[Empirical representation and convergence]
		\label{thm:empirical-cluster-counts}
		Let $\{X_t:t\in\Z^d\}$ be stationary and ergodic, let
		Assumption~\ref{ass:finite-cluster} hold, and suppose
		$\rho_u>0$. Then, for fixed $k\geq1$,
		\begin{align}
			\frac{N_k^*(\Lambda_n)}{|\Lambda_n|}
			&\overset{\rm a.s.}{\longrightarrow}w_k,
			&
			\frac{N^*(\Lambda_n)}{|\Lambda_n|}
			&\overset{\rm a.s.}{\longrightarrow}\rho_u,
			\label{eq:ideal-count-limits}\\
			\widetilde w_{k,n}
			:=\frac{N_k^\circ(\Lambda_n)}{|\Lambda_n|}
			&\overset{\rm a.s.}{\longrightarrow}w_k,
			&
			\widetilde\rho_{u,n}
			:=\frac{N^\circ(\Lambda_n)}{|\Lambda_n|}
			&\overset{\rm a.s.}{\longrightarrow}\rho_u.
			\label{eq:w-empirical}
		\end{align}
		Moreover, $N^\circ(\Lambda_n)>0$ eventually almost surely, and
		\begin{equation}\label{eq:empirical}
			\frac{N_k^\circ(\Lambda_n)}{N^\circ(\Lambda_n)}
			=\frac{\widetilde w_{k,n}}{\widetilde\rho_{u,n}}
			\overset{\rm a.s.}{\longrightarrow}
			\frac{w_k}{\rho_u}
			=\P(S_u=k).
		\end{equation}
		In fact, the whole empirical probability mass function is consistent in
		total variation:
		\begin{equation}\label{eq:empirical-total-variation}
			\sum_{k=1}^{\infty}
			\left|
			\frac{N_k^\circ(\Lambda_n)}{N^\circ(\Lambda_n)}
			-\P(S_u=k)
			\right|
			\overset{\rm a.s.}{\longrightarrow}0.
		\end{equation}
		For every finite $\Lambda$, the ideal counts also satisfy
		\begin{equation}\label{eq:ideal-count-expectations}
			\E N_k^*(\Lambda)=|\Lambda|w_k,
			\qquad
			\E N^*(\Lambda)=|\Lambda|\rho_u.
		\end{equation}
	\end{theorem}
	
	\begin{proof}
		The lexicographic order is invariant under translations. Consequently,
		$Y_{t,k}$ and $Y_t$ are translation-equivariant measurable functions of the
		field. Because the field is stationary and ergodic, both indicator fields are
		stationary and ergodic. Furthermore,
		\[
		\E Y_{o,k}
		=\P(R_u(o),|C_u(o)|=k)=w_k,
		\qquad
		\E Y_o=\P(R_u(o))=\rho_u.
		\]
		The pointwise ergodic theorem for the rectangular windows $\Lambda_n$ gives
		the two limits in \eqref{eq:ideal-count-limits}. Stationarity and linearity of
		expectation also give \eqref{eq:ideal-count-expectations}.
		
		It remains to show that discarding boundary affected clusters changes these
		counts by $o(|\Lambda_n|)$. On $R_u(t)$, define the closure radius
		\begin{equation}\label{eq:cluster-closure-radius}
			L_t
			=\max\{d_G(t,s):s\in C_u(t)\cup\mathcal N(C_u(t))\}.
		\end{equation}
		The graph is locally finite, and
		Assumption~\ref{ass:finite-cluster} implies that $L_t<\infty$ almost surely.
		
		Fix an integer $r\geq1$. If $t$ is a cluster root, $L_t\leq r$, and
		$t\notin\partial_r\Lambda_n$, then
		\[
		C_u(t)\cup\mathcal N(C_u(t))\subset\Lambda_n.
		\]
		That cluster is therefore included in
		$ \mathscr C_u^\circ(\Lambda_n)$. Conversely, every component in
		$ \mathscr C_u^\circ(\Lambda_n)$ is a full-lattice cluster with its root in
		$\Lambda_n$. Hence, for every $k$,
		\begin{align}
			0&\leq N_k^*(\Lambda_n)-N_k^\circ(\Lambda_n)
			\leq |\partial_r\Lambda_n|
			+\sum_{t\in\Lambda_n}Y_t\mathbbm{1}\{L_t>r\},
			\label{eq:size-count-boundary-error}\\
			0&\leq N^*(\Lambda_n)-N^\circ(\Lambda_n)
			\leq |\partial_r\Lambda_n|
			+\sum_{t\in\Lambda_n}Y_t\mathbbm{1}\{L_t>r\}.
			\label{eq:total-count-boundary-error}
		\end{align}
		Divide by $|\Lambda_n|$. The first term on the right tends to zero by
		\eqref{eq:empirical-boundary-negligible}. For fixed $r$, the pointwise
		ergodic theorem gives
		\[
		\frac1{|\Lambda_n|}\sum_{t\in\Lambda_n}
		Y_t\mathbbm{1}\{L_t>r\}
		\overset{\rm a.s.}{\longrightarrow}
		\P(R_u(o),L_o>r).
		\]
		Taking the limit superior in
		\eqref{eq:size-count-boundary-error}--\eqref{eq:total-count-boundary-error}
		and letting $r\to\infty$ yields zero, because $L_o<\infty$ almost surely
		on $R_u(o)$. Thus the observable and ideal counts have the same 
		limits, proving \eqref{eq:w-empirical}.
		
		Finally, $\widetilde\rho_{u,n}\to\rho_u>0$, so its denominator is eventually
		positive almost surely. Taking the ratio in \eqref{eq:w-empirical} and using
		Theorem~\ref{thm:CSD_d} proves \eqref{eq:empirical}. Because there are only
		countably many values of $k$, the convergence in \eqref{eq:empirical} holds
		simultaneously for all $k$ on a single event with probability one. On that event,
		both
		\[
		\left\{
		\frac{N_k^\circ(\Lambda_n)}{N^\circ(\Lambda_n)}:k\geq1
		\right\}
		\quad\text{and}\quad
		\{\P(S_u=k):k\geq1\}
		\]
		are probability mass functions. Thus Scheff\'e's
		lemma for counting measure yields
		\eqref{eq:empirical-total-variation}.
	\end{proof}
	
	The convergence in \eqref{eq:empirical}, together with \eqref{eq:observable-cluster-counts}, indicates that, as the window grows, the relative frequency of clusters of size $k$ converges to the theoretical cluster size probability in Theorem~\ref{thm:CSD_d}. This verifies that our theoretical formula agrees with the empirical
	cluster size distribution computed in simulations.
	
	\begin{comment}
	The theorem separates the empirical interpretation into an exact statement
	and an observable asymptotic statement.  If the full field were available,
	the expected number of size-$k$ roots per lattice site would be exactly
	$w_k$ by \eqref{eq:ideal-count-expectations}.  In a finite data window one
	cannot determine clusters that meet the unobserved exterior, but the
	interior counts in \eqref{eq:observable-cluster-counts} remove precisely
	that ambiguity.  The boundary argument above proves that this removal has
	vanishing effect per lattice site, without imposing a deterministic bound on
	cluster diameters.
	
	Here, we explain intuitively how Theorem~\ref{thm:CSD_d} agrees with the empirical
	cluster size distribution computed in simulations.  Suppose that the field is stationary and ergodic and is observed on expanding windows $\Lambda_n$; for simplicity, assume $|\Lambda_n|=n^d$.  To avoid ambiguity from clusters cut by the edge of the observation window, count only the interior
	components in $\mathscr C_u^\circ(\Lambda_n)$.  In simulations, the empirical probability mass function of size-$k$ clusters is the empirical proportion of size-$k$ clusters among all observed clusters, which is given by
	\begin{equation}\label{eq:empirical_2}
		\begin{split}
			\widetilde p_{k,n}
			&:=
			\frac{
				\#\{\text{interior size-$k$ clusters in $\Lambda_n$}\}}
			{\#\{\text{all interior clusters in $\Lambda_n$}\}}=
			\frac{N_k^\circ(\Lambda_n)}{N^\circ(\Lambda_n)}
			=
			\frac{N_k^\circ(\Lambda_n)/n^d}
			{N^\circ(\Lambda_n)/n^d}\\
			&\overset{\rm a.s.}{\longrightarrow}
			\frac{w_k}{\rho_u}
			=
			\frac{w_k}{\sum_{j=1}^{\infty}w_j}
			=
			\P(S_u=k),
			\qquad n\to\infty,
		\end{split}
	\end{equation}
	where the convergence is due to Theorem~\ref{thm:empirical-cluster-counts}. Thus, as the window grows, the
	relative frequency of clusters of size $k$ converges to the theoretical
	cluster size probability in Theorem~\ref{thm:CSD_d}.  This conclusion does
	not require cluster observations at different sites to be independent;
	stationarity and ergodicity provide the required spatial averaging.
	\end{comment}
	
	\subsection{Examples}
	
	\subsubsection{The one-dimensional specialization}
	
	Let $d=1$.  The nearest and Moore neighborhoods coincide, and the
	lexicographic order is the usual order on $\Z$.  As shown in
	\eqref{eq:root-shapes-1d}, $D_k=\{0,1,\ldots,k-1\}$ is the unique rooted
	connected shape of size $k$, and its exterior neighbor set is
	$\{-1,k\}$.  Therefore, the general weight \eqref{eq:w_d} reduces to
	\begin{equation}\label{eq:w_1D}
		\begin{split}
			w_k
			&=\P\bigl(X_{-1}\le u,\ X_0>u,\ X_1>u,\ldots,
			X_{k-1}>u,\ X_k\le u\bigr).
		\end{split}
	\end{equation}
	Moreover, the root event is simply
	$R_u(o)=\{X_{-1}\le u,X_0>u\}$.  Hence
	\begin{equation}\label{eq:w_1D_sum}
		\sum_{k=1}^{\infty}w_k
		=\P(X_{-1}\le u,X_0>u).
	\end{equation}
	
	
	Theorem~\ref{thm:CSD_d} now immediately gives the following one-dimensional
	formula.
	
	\begin{corollary}[One-dimensional specialization]\label{cor:CSD_1D}
		Let $\{X_t:t\in\Z\}$ be stationary, and let $u\in\R$ satisfy
		Assumption~\ref{ass:finite-cluster}.  If
		$\P(X_{-1}\le u,X_0>u)>0$, then, for $k=1,2,\ldots$,
		\begin{equation}\label{eq:CSD_1D}
			\P(S_u=k)
			=
			\frac{
				\P\bigl(X_{-1}\le u,\ X_0>u,\ X_1>u,\ldots,
				X_{k-1}>u,\ X_k\le u\bigr)}
			{\P(X_{-1}\le u,X_0>u)}.
		\end{equation}
	\end{corollary}
	
	\begin{proof}
		Substitute \eqref{eq:w_1D} and \eqref{eq:w_1D_sum} into the general formula
		\eqref{eq:CSD_d}.
	\end{proof}
	
	Thus, in one dimension the cluster size is the number of consecutive
	exceedances beginning at a transition from below to above the threshold.
	Figure~\ref{fig:simul-1D} provides a visual illustration, while
	\eqref{eq:empirical} gives the corresponding empirical relative frequency interpretation.
	
	\begin{example}[One-dimensional white noise]\label{example:WN_1D}
		Let $X_t$, $t\in\Z$, be an independent and identically distributed process
		with common cumulative distribution function $F$.  Set
		$p=F(u)$ and $q=1-p$, and assume $0<p<1$.  By
		\eqref{eq:CSD_1D},
		\begin{equation}\label{eq:WN_1D}
			\P(S_u=k)=\frac{p^2q^k}{pq}=pq^{k-1},
			\qquad k=1,2,\ldots.
		\end{equation}
		Hence the cluster size has a geometric distribution with success parameter
		$p$.
	\end{example}
	
	\subsubsection{A two-dimensional illustration}
	
	To illustrate the shape sum in \eqref{eq:w_d}, consider nearest-neighbor
	connectivity on $\Z^2$.  The rooted collections for $k\leq3$ are displayed
	immediately after the general definition in \eqref{eq:root3}.  They show that
	even for small $k$ several distinct rooted shapes may contribute to $w_k$;
	under Moore connectivity, diagonal adjacency produces still more shapes.
	
	Figure~\ref{fig:simul-2D} shows clusters and empirical cluster size
	distributions for a stationary Gaussian field on $\Z^2$ under both
	connectivity rules.  Because Moore connectivity includes diagonal sites, it
	generally produces larger connected components than nearest-neighbor
	connectivity.
	
	
	\begin{example}[Two-dimensional white noise]\label{example:WNF}
		Let $X_t$, $t\in\Z^2$, be an independent and identically distributed field
		with common CDF $F$, and set $p=F(u)$ and $q=1-p$.  Inspecting the rooted
		shapes in \eqref{eq:root3} and their exterior boundaries gives, under
		nearest-neighbor connectivity,
		\[
		w_1=p^4q,\qquad
		w_2=2p^6q^2,\qquad
		w_3=(2p^8+4p^7)q^3.
		\]
		Under Moore connectivity,
		\[
		w_1=p^8q,\qquad
		w_2=(2p^{10}+2p^{12})q^2,\qquad
		w_3=(6p^{12}+8p^{14}+4p^{15}+2p^{16})q^3.
		\]
		In general, the combinatorial complexity of
		$\mathcal{C}_k^{\rm root}$ grows rapidly with $k$.  Thus, closed-form
		expressions for $w_k$ are rarely available, although $w_k$ can be evaluated
		by enumerating shapes for moderate $k$ or approximated by Monte Carlo.
	\end{example}
	
	\begin{remark}[Computation of $w_k$]
		For large $k$, the probabilities $w_k$ in \eqref{eq:w_d} are difficult to
		compute analytically, and the normalizing constant
		$\sum_{j\ge1}w_j$ in \eqref{eq:CSD_d} is rarely available in closed form.
		The Monte Carlo estimators in Section \ref{sec:MC} provide alternatives.
		
		If the random field is isotropic, symmetry can substantially reduce the
		computational burden.  The probability
		$\P(X_D>u,X_{\mathcal{N}(D)}\le u)$ is the same for shapes related by a
		lattice rotation or reflection, so one may compute a single representative
		probability and multiply by its multiplicity.  For example, under
		nearest-neighbor connectivity on $\Z^2$, the two elements of
		$\mathcal{C}_2^{\rm root}$ are two orientations of the same geometric
		shape, while the six elements of $\mathcal{C}_3^{\rm root}$ in
		\eqref{eq:root3} reduce to two geometric types: straight triominoes and
		$L$-triominoes.
	\end{remark}
	
	\subsection{Equivalent representations and site-containing clusters}
	
	\subsubsection{A boundary representation of the normalizing constant}
	
	The general formulation also explains why the normalizing constant is simple
	in one dimension but generally nonlocal in higher dimensions.  Define
	\[
	B^-=\{t\in\mathcal{N}(o):t<_{\rm lex}o\},
	\qquad
	H^-=\{t\in\Z^d\setminus \mathcal{N}(o):t<_{\rm lex}o\}.
	\]
	For nearest-neighbor connectivity, $B^-=\{-e_i:i=1,\ldots,d\}$.  
	
	The origin is a cluster root precisely when it is an exceedance, its neighbors
	in $B^-$ are not exceedances, and there is no exceedance path from $o$ to a
	site in $H^-$.  Consequently,
	\begin{equation}\label{eq:sum-w}
		\begin{split}
			\rho_u
			&=\P(X_o>u,X_{B^-}\le u)\\
			&\quad-
			\P\bigl(X_o>u,X_{B^-}\le u,
			\text{and there exists an exceedance path from $o$ to $H^-$}\bigr).
		\end{split}
	\end{equation}
	When $d=1$, the site $-1$ separates $o=0$ from every site in $H^-$, so the
	last probability in \eqref{eq:sum-w} vanishes.  This recovers
	\eqref{eq:w_1D_sum}.  When $d\ge2$, an excursion path can go around the lower
	neighbors in $B^-$, and the correction term is generally difficult to
	evaluate.
	
	\subsubsection{Clusters containing the origin versus rooted at the origin}
	
	For $k\ge1$, let
	\[
	\mathcal{C}_k^{\rm inside}
	=\{D\subset\Z^d:|D|=k,\ D\text{ is connected, and }o\in D\}.
	\]
	Here the origin may occupy any one of the $k$ sites of the cluster.  Translating
	each rooted shape so that each of its sites is placed at $o$ gives
	\begin{equation}\label{eq:C-inside}
		|\mathcal{C}_k^{\rm inside}|=k|\mathcal{C}_k^{\rm root}|.
	\end{equation}
	Let
	\begin{equation}\label{eq:w-inside}
		\begin{split}
			w_k^{\rm inside}
			&:=\P\bigl(X_o>u,\ |C_u(o)|=k\bigr)\\
			&=\sum_{D\in\mathcal{C}_k^{\rm inside}}
			\P\bigl(X_D>u,X_{\mathcal{N}(D)}\le u\bigr).
		\end{split}
	\end{equation}
	By stationarity, each cluster of size $k$ contributes once to the rooted count
	and $k$ times to the site-containing count.  Hence
	\begin{equation}\label{eq:w-inside-w}
		w_k^{\rm inside}=kw_k.
	\end{equation}
	
	\begin{corollary}\label{cor:CSD_d}
		Under the assumptions of Theorem~\ref{thm:CSD_d},
		\[
		\P(S_u=k)
		=
		\frac{w_k^{\rm inside}/k}
		{\sum_{j=1}^{\infty}w_j^{\rm inside}/j},
		\qquad k=1,2,\ldots,
		\]
		where $w_k^{\rm inside}$ is defined in \eqref{eq:w-inside}.
	\end{corollary}
	
	The relation \eqref{eq:w-inside-w} also leads to unbiased origin-based and
	spatial estimators of $w_k$. We develop their exact algorithms,
	finite-sample unbiasedness, control of the omitted normalization tail, and
	estimation of the spatial Monte Carlo variance in Section~\ref{sec:MC}.
	
	\section{The peak-based cluster size distribution}\label{sec:peak-based-cluster}
	When the random field $X_t$ is nonstationary, the cluster size distribution in \eqref{eq:CSD_def} becomes challenging to analyze. This difficulty arises because, in the nonstationary setting, the conditional probability in \eqref{eq:CSD_def} depends on the spatial location of the observed cluster $C_u$, as the covariance structure of the field varies across the domain. Consequently, there is generally no canonical translation-invariant cluster size distribution, since the law of a cluster depends on its spatial location. A more appropriate formulation is a local, spatially varying cluster size distribution that reflects the heterogeneity in the covariance structure of the field. 
	
	A second motivation comes from the fact that every cluster contains at least one local maximum (typically only a few), which can serve as a natural representative of the cluster. Moreover, for high thresholds $u$, clusters tend to be dominated by a single prominent peak. Since peak analysis is itself a mature and active area \cite{Cheng:2020,CS15,CS18,CS17,Fondeville:2018,Huser:2014}, it is natural and valuable to study the peak-based cluster size distribution from both theoretical and applied perspectives. Specifically, given that a local maximum exceeding $u$ is observed at $t$, we consider the distribution of the size of the cluster containing that peak. 
	
		\begin{definition}\label{def:peak-cluster-size} For a discrete random field $\{X_t: t\in \Z^d\}$, the distribution of the peak-based cluster size at $t\in \Z^d$, denoted $S_u^{\rm peak}(t)$, is defined as 
		\begin{equation}\label{eq:CSD_peak_1}
			\begin{split}
				\P(S_u^{\rm peak}(t)=k) = \P( |C_u(t)|=k \ | \ t \ \text{\rm is a local maximum and } X_t>u), \quad k\ge 1,
			\end{split}
		\end{equation}
		provided the conditional event has a positive probability. 
	\end{definition}
	
	The peak-based formulation is inherently local: unlike the rooted cluster distribution, its normalization depends only on the finite-dimensional distribution of the field in a neighborhood of the peak location. This locality makes the peak-based approach particularly suitable for nonstationary random fields.
	
	In signal detection problems, one is often interested in peaks that are not only high in amplitude but also substantial in spatial extent. This is particularly relevant in neuroimaging applications, where extended activation patterns are more scientifically interpretable than isolated spikes \cite{Chumbley:2010,Poline:1997,Worsley:1996a,Zhang:2009}. A peak-based formulation of the cluster size distribution provides a natural and efficient framework for quantifying such effects and assessing their significance, through $p$-values, when detecting spatially extended signals.
	
	
	\subsection{General $d$-dimensional formulation}
	Let $\{X_t:t\in\Z^d\}$ be a real-valued random field and fix $u\in\R$. For $t\in\Z^d$, define the peak-over-threshold event
	\begin{equation}\label{eq:peak-event}
		A_u^{\rm peak}(t)
		:=
		\left\{X_t>u,\
		X_t\ge \max_{s\in\mathcal N(t)}X_s\right\},
	\end{equation}
	and write
	\begin{equation}\label{eq:peak-intensity-local}
		\rho_u^{\rm peak}(t):=\P\bigl(A_u^{\rm peak}(t)\bigr).
	\end{equation}
	Thus, $A_u^{\rm peak}(t)$ occurs precisely when $t$ is a local maximum
	and exceeds $u$. If the field is stationary, then $\rho_u^{\rm peak}(t)$ does not depend on $t$, we simply write it as $\rho_u^{\rm peak}$.
	
	Suppose Assumption~\ref{ass:finite-cluster} holds and
	$\rho_u^{\rm peak}(t)>0$.  Then the peak-based cluster size distribution
		at $t$ in \eqref{eq:CSD_peak_1} can be written as
	\begin{equation}\label{eq:CSD_peak}
		\P\bigl(S_u^{\rm peak}(t)=k\bigr) =
		\P\bigl(|C_u(t)|=k\mid A_u^{\rm peak}(t)\bigr),
		\qquad k=1,2,\ldots.
	\end{equation}
	For a nonstationary field, this distribution may depend on $t$.  For a
	stationary field it does not, and we then write $S_u^{\rm peak}$ instead of
	$S_u^{\rm peak}(t)$.
	
	
	For $k\geq1$, define the deterministic collection of connected shapes of size
	$k$ containing the distinguished site $t$ by
	\begin{equation}\label{eq:C-peak}
		\mathcal C_k^{\rm peak}(t)
		:=
		\left\{
		D\subset\Z^d:
		|D|=k,\ D\text{ is connected, and }t\in D
		\right\}.
	\end{equation}
	The superscript ``peak'' records that $t$ is the candidate peak location; the
	shapes themselves are deterministic and do not depend on the field values.
	Translation gives
	\begin{equation}\label{eq:C-peak-cardinality}
		|\mathcal C_k^{\rm peak}(t)|
		=|\mathcal C_k^{\rm inside}|
		=k|\mathcal C_k^{\rm root}|.
	\end{equation}
	
	For example, in $\Z^2$ with nearest-neighbor connectivity,
	\begin{equation*}
		\begin{split}
			\mathcal C_1^{\rm peak}(t) &= \bigl\{\{t\}\bigr\}, \\
			\mathcal C_2^{\rm peak}(t) &= \bigl\{\{t,t+e_1\}, \{t,t-e_1\},
			\{t,t+e_2\}, \{t,t-e_2\}\bigr\}.
		\end{split}
	\end{equation*}
	There are $18$ shapes in $\mathcal C_3^{\rm peak}(t)$.  Under Moore
	connectivity, diagonal shapes are also present; in particular,
	$\mathcal C_2^{\rm peak}(t)$ contains four additional diagonal pairs.
	
	For $D\in\mathcal C_k^{\rm peak}(t)$, the event that $D$ is exactly the
	excursion cluster containing $t$ and that $t$ is a peak above $u$ is
	\begin{equation}\label{eq:peak-shape-event}
		\{C_u(t)=D\}\cap A_u^{\rm peak}(t)
		=
		\left\{
		X_D>u,\ X_{\mathcal N(D)}\leq u,\
		X_t\ge \max_{s\in\mathcal N(t)}X_s
		\right\}.
	\end{equation}
	Define the size-$k$ peak intensity at $t$ by
	\begin{equation}\label{eq:w-peak}
		\begin{split}
			w_k^{\rm peak}(t)
			&:=\sum_{D\in\mathcal C_k^{\rm peak}(t)}
			\P\left(
			X_D>u,\ X_{\mathcal N(D)}\leq u,\
			X_t\ge \max_{s\in\mathcal N(t)}X_s
			\right).
		\end{split}
	\end{equation}
	Specifically, $w_k^{\rm peak}(t)$ is the probability that the specified site $t$ is
	a peak above $u$ and that its associated excursion cluster has size $k$. If $X$ is stationary, then $w_k^{\rm peak}(t)$ does not depend on $t$, and we simply write it as $w_k^{\rm peak}$.
	
\begin{theorem}[Peak-based cluster size formula]
	\label{thm:peak-CSD}
	Let $\{X_t:t\in\Z^d\}$ be a random field, not necessarily stationary, and
	let $u\in\R$ satisfy Assumption~\ref{ass:finite-cluster}.  Fix
	$t\in\Z^d$ and suppose $\rho_u^{\rm peak}(t)$ defined in \eqref{eq:peak-intensity-local} is positive. Then, for every $k\geq1$,
	\begin{equation}\label{eq:peak-CSD}
		\P\bigl(S_u^{\rm peak}(t)=k\bigr)
		=
		\frac{w_k^{\rm peak}(t)}
		{\rho_u^{\rm peak}(t)}, \qquad \rho_u^{\rm peak}(t)
		=\sum_{j=1}^{\infty}w_j^{\rm peak}(t).
	\end{equation}
	where $w_k^{\rm peak}(t)$ is given by \eqref{eq:w-peak}.  
\end{theorem}

\begin{proof}
	By expression \eqref{eq:CSD_peak}, we obtain
	\begin{equation}\label{eq:CSD_peak_ratio}
		\P\bigl(S_u^{\rm peak}(t)=k\bigr) = \frac{\P\bigl(A_u^{\rm peak}(t),\ |C_u(t)|=k\bigr)}{\P(A_u^{\rm peak}(t))},
		\qquad k=1,2,\ldots.
	\end{equation}
	For a fixed $k$, the events on the left hand side of
	\eqref{eq:peak-shape-event}, indexed by
	$D\in\mathcal C_k^{\rm peak}(t)$, are pairwise disjoint.  Their union over all possible $D\in\mathcal C_k^{\rm peak}(t)$ is
	exactly
	\[
	A_u^{\rm peak}(t)\cap\{|C_u(t)|=k\}.
	\]
	Taking probabilities gives 
	\[
	w_k^{\rm peak}(t) =\P\bigl(A_u^{\rm peak}(t),\ |C_u(t)|=k\bigr).
	\]
	Dividing by
	$\P(A_u^{\rm peak}(t))=\rho_u^{\rm peak}(t)$ and using
	\eqref{eq:CSD_peak_ratio} proves the first equality in \eqref{eq:peak-CSD}.
	
	By Assumption~\ref{ass:finite-cluster}, the cluster $C_u(t)$ has a finite
	size a.s. on $A_u^{\rm peak}(t)$.  Consequently,
	\[
	A_u^{\rm peak}(t)
	=
	\bigsqcup_{j=1}^{\infty}
	\left(A_u^{\rm peak}(t)\cap\{|C_u(t)|=j\}\right).
	\]
	Taking probabilities proves the second equality in \eqref{eq:peak-CSD}.
\end{proof}
	
	Theorem~\ref{thm:peak-CSD} is local in its normalization: the denominator
	$\rho_u^{\rm peak}(t)$ involves only the joint distribution of
	$\{X_s:s\in\{t\}\cup\mathcal N(t)\}$.  This is the principal advantage over
	the rooted cluster distribution, whose normalizing constant generally
	contains a global connectivity correction.  The numerator in
	\eqref{eq:w-peak}, however, still becomes combinatorially demanding as $k$
	increases.
	
	\subsection{Examples}
	
	Consider first the one-dimensional case $d=1$.  If a size-$k$ cluster
	contains the specified peak site $t$, let $j\in\{0,\ldots,k-1\}$ be the
	number of cluster sites strictly to the right of $t$.  The corresponding
	interval is
	\begin{equation*}
		D_{k,j}(t)
		:=\{t-k+j+1,\ldots,t+j\},
		\qquad j=0,\ldots,k-1,
	\end{equation*}
	and
	\[
	\mathcal C_k^{\rm peak}(t)
	=\{D_{k,j}(t):j=0,\ldots,k-1\}.
	\]
	The two exterior neighbors of $D_{k,j}(t)$ are $t-k+j$ and $t+j+1$.
	Substitution into the general formula \eqref{eq:w-peak} gives
	\begin{equation}\label{eq:peak-GP}
		\begin{split}
			\P(S_u^{\rm peak}(t)=k)
			={1\over\rho_u^{\rm peak}(t)}
			\sum_{j=0}^{k-1}\P\bigl(&X_{t-k+j}\leq u,
			X_{t-k+j+1}>u,\ldots,X_{t+j}>u,\\[-1mm]
			&X_{t+j+1}\leq u,
			X_t\ge \max\{X_{t-1},X_{t+1}\}\bigr),
		\end{split}
	\end{equation}
	where
	\[
	\rho_u^{\rm peak}(t) = \P\left(X_t>u, \, X_t\ge \max\{X_{t-1},X_{t+1}\}\right).
	\]
	In particular, for $k=1$,
	\begin{equation}\label{eq:peak1}
		\P(S_u^{\rm peak}(t)=1)
		=
		\P\bigl(X_{t-1}\leq u,X_{t+1}\leq u
		\mid A_u^{\rm peak}(t)\bigr).
	\end{equation}
	For $k=2$, the cluster lies either to the right or to the left of $t$, so
	\begin{equation}\label{eq:peak2}
		\begin{split}
			\P(S_u^{\rm peak}(t)=2)
			={1\over\rho_u^{\rm peak}(t)}\bigl[&
			\P(X_{t-1}\leq u,X_t\ge X_{t+1}>u,X_{t+2}\leq u)\\
			&+\P(X_{t-2}\leq u,X_t\ge X_{t-1}>u,X_{t+1}\leq u)
			\bigr].
		\end{split}
	\end{equation}
	If the law of the process is invariant under reflection about $t$, the terms
	in \eqref{eq:peak-GP} indexed by $j$ and $k-1-j$ coincide.  
	
	\begin{example}[One-dimensional white noise]
		\label{example:peak-GWN} 
		Let $X_t$, $t\in \Z$, be a white noise process with common continuous CDF $F$. Let $p=F(u)$ and $q=1-p$. Note that, if $Y_1, Y_2$ and $Y_3$ are i.i.d. with CDF $F$ and density $f$, then 
		\[
		\P(Y_1\ge Y_2, Y_1\ge Y_3 \, | \, Y_1=x) = \P(Y_2\le x)^2 = F(x)^2, 
		\]
		so
		\[
		\rho_u^{\rm peak} = \P(Y_1>u, Y_1\ge Y_2, Y_1\ge Y_3) = \int_u^\infty f(x)F(x)^2dx = \frac{1-p^3}{3}.
		\]
		Similarly,
		\begin{equation*}
			\begin{split}
				\P(Y_1\ge Y_2>u) = \frac{q^2}{2},\quad
				\P(Y_1\ge Y_2>u, Y_1\ge Y_3>u) =\frac{q^3}{3}.
			\end{split}
		\end{equation*}
		It follows from \eqref{eq:peak1} and \eqref{eq:peak2} that
		\begin{equation*}
			\begin{split}
				\P(S_u^{\rm peak}=1) = \frac{3p^2q}{1-p^3},\quad
				\P(S_u^{\rm peak}=2) = \frac{3p^2q^2}{1-p^3}.
			\end{split}
		\end{equation*}
		Similarly, 
		\begin{equation*}
			\begin{split}
				\P(S_u^{\rm peak}=3) &= \frac{2\P(X_{t-1}\le u, X_t\ge X_{t+1}>u, X_{t+2}>u, X_{t+3}\le u))}{\rho_u^{\rm peak}}\\
				&\quad + \frac{\P(X_{t-2}\le u, X_t\ge X_{t-1}> u, X_t\ge X_{t+1}>u, X_{t+2}\le u)}{\rho_u^{\rm peak}}\\
				&= \frac{p^2q^3 + p^2q^3/3}{(1-p^3)/3} = \frac{4p^2q^3}{1-p^3}.
			\end{split}
		\end{equation*}
		In general, we have
		\begin{equation}\label{eq:peak-GWN-1D}
			\P(S_u^{\rm peak}=k) = \begin{cases}
				\frac{3p^2q}{1-p^3}, \quad  &k=1,\\
				\frac{(k+1)p^2q^k}{1-p^3}, \quad  &k\ge 2.
			\end{cases}
		\end{equation}
		%\begin{equation*}
		%	\begin{split}
			%		\P(S_u^{\rm peak}(t)=4) = \frac{5p^2q^4}{1-p^3}, \quad \P(S_u^{\rm %peak}(t)=5) = \frac{6p^2q^5}{1-p^3}, \quad \P(S_u^{\rm peak}(t)=6) = %\frac{7p^2q^6}{1-p^3}.
			%	\end{split}
		%\end{equation*}
	\end{example}
	
	
	
	\begin{example}[Two-dimensional white noise]
		\label{example:WNF-peak}
		Let $\{X_t:t\in\Z^2\}$ be i.i.d. with a continuous CDF $F$, and set
		$p=F(u)$ and $q=1-p$.  Under nearest-neighbor connectivity, a site is
		compared with four neighbors, and hence
		\[
		\rho_u^{\rm peak}
		=\int_u^\infty F(x)^4\,dF(x)
		=\frac{1-p^5}{5}.
		\]
		Enumerating the shapes in $\mathcal C_k^{\rm peak}(t)$ and distinguishing
		whether the candidate peak has one or two neighbors inside the cluster
		gives
		\begin{equation*}
			\begin{split}
				w_1^{\rm peak} = p^4q, \quad w_2^{\rm peak} = 2p^6q^2, \quad w_3^{\rm peak} = \frac{8}{3}(2p^7+p^8)q^3.
			\end{split}
		\end{equation*}
		Taking the ratio yields the following probabilities:
		\begin{equation*}
			\begin{split}
				\P(S_u^{\rm peak}=1) = \frac{5p^4q}{1-p^5}, \quad \P(S_u^{\rm peak}=2) = \frac{10p^6q^2}{1-p^5}, \quad 
				\P(S_u^{\rm peak}=3) = \frac{40(2p^7+p^8)q^3}{3(1-p^5)}.
			\end{split}
		\end{equation*}
		
		Under Moore connectivity, the site is compared with eight neighbors, so
		$\rho_u^{\rm peak}=(1-p^9)/9$, and
		\begin{equation*}
			\begin{split}
				w_1^{\rm peak} &= p^8q, \quad w_2^{\rm peak}(t) = 2(p^{10}+p^{12})q^2,\\
				w_3^{\rm peak} &= \frac{4}{3}(5p^{12}+8p^{14}+4p^{15}+2p^{16})q^3.
			\end{split}
		\end{equation*}
		Consequently,
		\begin{equation*}
			\begin{split}
				\P(S_u^{\rm peak}=1) &= \frac{9p^8q}{1-p^9}, \quad \P(S_u^{\rm peak}=2) = \frac{18(p^{10}+p^{12})q^2}{1-p^9}, \\
				\P(S_u^{\rm peak}=3) &= \frac{12(5p^{12}+8p^{14}+4p^{15}+2p^{16})q^3}{1-p^9}.
			\end{split}
		\end{equation*}
		As $k$ increases, the combinatorial complexity of $\mathcal{C}_k^{\rm peak}$ grows quickly, so it is difficult to obtain closed-form expressions for $w_k^{\rm peak}$. Nevertheless, $w_k^{\rm peak}$ can be evaluated numerically by enumerating cluster shapes or estimated by Monte Carlo in Section \ref{sec:MC}.
	\end{example}
	
\subsection{Empirical peak counts and consistency}
\label{sec:empirical-peak-counts}

We now assume that the field is stationary.  
For $t\in\Z^d$, define the peak indicators
\begin{equation}\label{eq:empirical-peak-indicators}
	Z_{t,k}
	=\mathbbm{1}\{A_u^{\rm peak}(t),|C_u(t)|=k\},
	\qquad
	Z_t=\mathbbm{1}\{A_u^{\rm peak}(t)\}.
\end{equation}
For a finite window $\Lambda$, define the ideal peak counts
\begin{equation}\label{eq:ideal-peak-counts}
	N_k^{{\rm peak},*}(\Lambda)
	=\sum_{t\in\Lambda}Z_{t,k},
	\qquad
	N^{{\rm peak},*}(\Lambda)
	=\sum_{t\in\Lambda}Z_t.
\end{equation}
Thus, $N_k^{{\rm peak},*}(\Lambda)$ counts the peak sites in $\Lambda$ whose
associated full-lattice clusters have size $k$, whereas
$N^{{\rm peak},*}(\Lambda)$ counts all peak sites above $u$ in $\Lambda$.
These are counts of \emph{peaks}, not counts of distinct clusters.  In
particular, if one excursion cluster contains several local maxima, it is
represented once for each of those maxima in the peak-based distribution.
Under Assumption~\ref{ass:finite-cluster}, $N^{{\rm peak},*}(\Lambda)
=\sum_{j=1}^{\infty}N_j^{{\rm peak},*}(\Lambda)$.

As in Section~\ref{sec:empirical-cluster-counts}, the superscript $*$ means
that a peak is counted according to its full-lattice cluster, even if that
cluster crosses the boundary of $\Lambda$.  The ideal counts may therefore
depend on unobserved values outside $\Lambda$.

To obtain observable counts, recall the interior cluster
collection $\mathscr C_u^\circ(\Lambda)$ from
\eqref{eq:observable-interior-components}.  For
$D\in\mathscr C_u^\circ(\Lambda)$, define its set of observed peak sites by
\begin{equation*}
	\operatorname{Peak}_u(D)
	:=
	\left\{
	t\in D:X_t\ge \max_{s\in\mathcal N(t)}X_s
	\right\}.
\end{equation*}
Every neighbor of every $t\in D$ is observed because
$D\cup\mathcal N(D)\subset\Lambda$.  Hence
$\operatorname{Peak}_u(D)$ and $|D|$ are both determined by the observations
inside $\Lambda$.  Define
\begin{equation}\label{eq:observable-peak-counts}
	\begin{split}
	N_k^{{\rm peak},\circ}(\Lambda)
	&=
	\sum_{D\in\mathscr C_u^\circ(\Lambda)}
	\mathbbm{1}\{|D|=k\}\,|\operatorname{Peak}_u(D)|\\
	&=\#\{\text{observed peaks with associated size-$k$ interior clusters in $\Lambda$}\};
	\\
	N^{{\rm peak},\circ}(\Lambda) 
	&=
	\sum_{D\in\mathscr C_u^\circ(\Lambda)}
	|\operatorname{Peak}_u(D)|\\
	&=\#\{\text{all observed peaks with associated interior clusters in $\Lambda$}\}.
	\end{split}
\end{equation}

Let $\Lambda_n$ be the expanding rectangular windows in
\eqref{eq:empirical-rectangular-windows}, so that their
boundaries satisfy \eqref{eq:empirical-boundary-negligible}. The following theorem follows directly from similar arguments for proving Theorem \ref{thm:empirical-cluster-counts}; the detailed proof is omitted here.

\begin{theorem}[Consistency of empirical peak frequencies]
	\label{thm:empirical-peak-counts}
	Let $\{X_t:t\in\Z^d\}$ be stationary and ergodic, let
	Assumption~\ref{ass:finite-cluster} hold, and suppose
	$\rho_u^{\rm peak}>0$.  Then, for fixed $k\geq1$,
	\begin{align}
		\frac{N_k^{{\rm peak},*}(\Lambda_n)}{|\Lambda_n|}
		&\overset{\rm a.s.}{\longrightarrow}w_k^{\rm peak},
		&
		\frac{N^{{\rm peak},*}(\Lambda_n)}{|\Lambda_n|}
		&\overset{\rm a.s.}{\longrightarrow}\rho_u^{\rm peak},
		\nonumber
		\\
		\widetilde w_{k,n}^{\rm peak}
		:=\frac{N_k^{{\rm peak},\circ}(\Lambda_n)}{|\Lambda_n|}
		&\overset{\rm a.s.}{\longrightarrow}w_k^{\rm peak},
		&
		\widetilde \rho_{u,n}^{\rm peak}
		:=\frac{N^{{\rm peak},\circ}(\Lambda_n)}{|\Lambda_n|}
		&\overset{\rm a.s.}{\longrightarrow}\rho_u^{\rm peak}.
		\label{eq:w-peak-empirical}
	\end{align}
	Moreover, $N^{{\rm peak},\circ}(\Lambda_n)>0$ eventually almost surely,
	and
	\begin{equation}\label{eq:empirical-peak-ratio}
		\frac{N_k^{{\rm peak},\circ}(\Lambda_n)}
		{N^{{\rm peak},\circ}(\Lambda_n)} = \frac{\widetilde w_{k,n}^{\rm peak}}{\widetilde \rho_{u,n}^{\rm peak}}
		\overset{\rm a.s.}{\longrightarrow}
		\frac{w_k^{\rm peak}}{\rho_u^{\rm peak}}
		=\P(S_u^{\rm peak}=k).
	\end{equation}
	In fact, the whole empirical probability mass function is consistent in total variation:
	\begin{equation*}
		\sum_{k=1}^{\infty}
		\left|
		\frac{N_k^{{\rm peak},\circ}(\Lambda_n)}
		{N^{{\rm peak},\circ}(\Lambda_n)}
		-\P(S_u^{\rm peak}=k)
		\right|
		\overset{\rm a.s.}{\longrightarrow}0.
	\end{equation*}
	For every finite $\Lambda$, the ideal counts satisfy the exact expectation
	identities
	\begin{equation*}
		\E N_k^{{\rm peak},*}(\Lambda)=|\Lambda|w_k^{\rm peak},
		\qquad
		\E N^{{\rm peak},*}(\Lambda)=|\Lambda|\rho_u^{\rm peak}.
	\end{equation*}
\end{theorem}

\begin{comment}
\begin{proof}
	The indicator fields $\{Z_{t,k}\}$ and $\{Z_t\}$ are
	translation-equivariant measurable functions of the original field.
	Stationarity and ergodicity of $X$ therefore imply stationarity and
	ergodicity of both indicator fields.  Their means are
	\[
	\E Z_{o,k}=w_k^{\rm peak},
	\qquad
	\E Z_o=\rho_u^{\rm peak}.
	\]
	The pointwise ergodic theorem gives
	\eqref{eq:ideal-peak-count-limits}, and stationarity together with
	linearity of expectation gives \eqref{eq:ideal-peak-count-expectations}.
	
	It remains to compare the ideal counts with the observable counts.  On
	$A_u^{\rm peak}(t)$, define the cluster-closure radius
	\begin{equation}\label{eq:peak-cluster-closure-radius}
		L_t^{\rm peak}
		:=
		\max\{d_G(t,s):s\in C_u(t)\cup\mathcal N(C_u(t))\}.
	\end{equation}
	The graph is
	locally finite, and Assumption~\ref{ass:finite-cluster} implies
	$L_t^{\rm peak}<\infty$ almost surely.
	
	Fix $r\geq1$.  If $t\in\Lambda_n$ is a peak above $u$,
	$t\notin\partial_r\Lambda_n$, and $L_t^{\rm peak}\leq r$, then
	\[
	C_u(t)\cup\mathcal N(C_u(t))\subset\Lambda_n.
	\]
	Thus, $C_u(t)$ belongs to $\mathscr C_u^\circ(\Lambda_n)$, and the peak
	$t$ is included in the corresponding observable count.  Conversely, every
	peak counted in \eqref{eq:observable-size-peak-count} or
	\eqref{eq:observable-total-peak-count} is a genuine full-lattice peak whose
	location belongs to $\Lambda_n$.  Therefore,
	\begin{align}
		0
		&\leq
		N_k^{{\rm peak},*}(\Lambda_n)
		-N_k^{{\rm peak},\circ}(\Lambda_n)\leq
		|\partial_r\Lambda_n|
		+\sum_{t\in\Lambda_n}
		Z_t\mathbbm{1}\{L_t^{\rm peak}>r\},
		\label{eq:size-peak-boundary-error}\\
		0
		&\leq
		N^{{\rm peak},*}(\Lambda_n)
		-N^{{\rm peak},\circ}(\Lambda_n)\leq
		|\partial_r\Lambda_n|
		+\sum_{t\in\Lambda_n}
		Z_t\mathbbm{1}\{L_t^{\rm peak}>r\}.
		\label{eq:total-peak-boundary-error}
	\end{align}
	After division by $|\Lambda_n|$, the boundary term tends to zero.  For
	each fixed integer $r$, another application of the pointwise ergodic
	theorem gives
	\[
	\frac1{|\Lambda_n|}\sum_{t\in\Lambda_n}
	Z_t\mathbbm{1}\{L_t^{\rm peak}>r\}
	\overset{\rm a.s.}{\to}
	\P\bigl(A_u^{\rm peak}(o),L_o^{\rm peak}>r\bigr).
	\]
	Taking limit superiors and then
	letting $r\to\infty$ gives zero because $L_o^{\rm peak}<\infty$ almost
	surely on $A_u^{\rm peak}(o)$.  Hence the observable and ideal peak counts
	have the same limits, yielding \eqref{eq:observable-peak-count-limits}.
	
	Since $\widetilde\pi_{u,n}^{\rm peak}\to\rho_u^{\rm peak}>0$, the observable
	denominator is eventually positive almost surely.  Taking the ratio and
	using Theorem~\ref{thm:peak-CSD} proves
	\eqref{eq:empirical-peak-ratio}.  The convergence holds simultaneously for
	all $k$ on a probability-one event. Scheff\'e's lemma for counting measure proves \eqref{eq:empirical-peak-total-variation}.
\end{proof}
\end{comment}

Combining \eqref{eq:observable-peak-counts} and the convergence in \eqref{eq:empirical-peak-ratio}, we see that, as the window grows, the relative frequency of peaks in size-$k$ clusters converges to the theoretical peak-based cluster size probability in Theorem~\ref{thm:peak-CSD}. This verifies that our theoretical formula agrees with the empirical
peak-based cluster size distribution computed in simulations.

For a nonstationary field, the analogous fixed-site result follows directly
from the strong law applied to independent field realizations. That is, the ratio of the number of peak events at $t$ with size-$k$ cluster to the total number of peak events at $t$ converges almost surely to
$\P(S_u^{\rm peak}(t)=k)$.


\begin{comment}
The fixed-site analogue remains available when the field is nonstationary.

\begin{corollary}[Independent simulations at a fixed site]
	\label{cor:replicated-peak-empirical}
	Let $X^{(1)},X^{(2)},\ldots$ be independent copies of a random field
	satisfying Assumption~\ref{ass:finite-cluster}.  Fix $t\in\Z^d$ and suppose
	$\rho_u^{\rm peak}(t)>0$.  Define
	\begin{equation}\label{eq:replicated-peak-indicators}
		B_{m,k}
		:=\mathbbm{1}\{A_{u,m}^{\rm peak}(t),|C_{u,m}(t)|=k\},
		\qquad
		B_m:=\mathbbm{1}\{A_{u,m}^{\rm peak}(t)\},
	\end{equation}
	where the subscript $m$ denotes quantities computed from $X^{(m)}$.
	Then
	\begin{equation}\label{eq:replicated-peak-limits}
		\frac1M\sum_{m=1}^M B_{m,k}
		\to w_k^{\rm peak}(t),
		\qquad
		\frac1M\sum_{m=1}^M B_m
		\to\rho_u^{\rm peak}(t), \quad {\rm a.s}.
	\end{equation}
	Moreover, $\sum_{m=1}^M B_m>0$ eventually almost surely, and
	\begin{equation}\label{eq:replicated-peak-ratio}
		\frac{\sum_{m=1}^M B_{m,k}}
		{\sum_{m=1}^M B_m}
		\to
		\P\bigl(S_u^{\rm peak}(t)=k\bigr), \quad {\rm a.s}.
	\end{equation}
\end{corollary}

\begin{proof}
	For fixed $k$, the pairs $(B_{m,k},B_m)$ are i.i.d. and bounded, with
	means $w_k^{\rm peak}(t)$ and $\rho_u^{\rm peak}(t)$, respectively.  The
	strong law of large numbers proves \eqref{eq:replicated-peak-limits}.  The
	second limit is strictly positive, so the denominator in
	\eqref{eq:replicated-peak-ratio} is eventually nonzero.  Taking the ratio
	and applying Theorem~\ref{thm:peak-CSD} proves the claim.
\end{proof}

This corollary formalizes the empirical distribution obtained by repeatedly
simulating a possibly nonstationary field and examining the same location
$t$.  If only a finite region is simulated, the cluster associated with $t$
must additionally be interior with respect to the domain, or the simulation domain must
contain a sufficient guard band.
\end{comment}
	
\section{Comparison between Exact and Peak-based Cluster Size Distributions for High Thresholds}\label{sec:high-u}

For many stationary fields, a high excursion cluster typically
contains a unique local maximum.  On realizations where every relevant
cluster contains exactly one peak, sampling clusters by their roots and
sampling them by their peaks give the same cluster weights.  This explains
the useful high-threshold approximation
\[
\P(S_u^{\rm peak}=k)\approx\P(S_u=k), \qquad \text{ for large $u$}.
\]
We will demonstrate rigorously this approximation under some mild conditions.

For a threshold $u$, define, on the root event $R_u(o)$,
\begin{equation}\label{eq:number-peaks-root-cluster}
	M_u(o)
	:=
	\sum_{t\in C_u(o)}
	\mathbbm{1}\{A_u^{\rm peak}(t)\},
\end{equation}
the number of local maxima in the excursion cluster rooted at the
origin. Since $C_u(o)$ is finite on $R_u(o)$, one has $M_u(o)\geq1$. 

\begin{theorem}[Root--peak comparison]
	\label{thm:root-peak-comparison}
	Suppose that the field is stationary, Assumption~\ref{ass:finite-cluster}
	holds, and $\rho_u=\P(R_u(o))>0$. Then
	\begin{equation}\label{eq:root-peak-TV-bound}
		\sum_{k=1}^{\infty} \left|
		\P(S_u^{\rm peak}=k)
		-\P(S_u=k)
		\right|
		\leq
		\frac{\E\left[M_u(o)-1\mid R_u(o)\right]}
		{\E\left[M_u(o)\mid R_u(o)\right]}
		=
		1-\frac{\rho_u}{\rho_u^{\rm peak}}.
	\end{equation}
	Hence the two distributions are identical when $M_u(o)=1$ almost surely on
	$R_u(o)$, and they are asymptotically equivalent if
	\[
	\E\left[M_u(o)-1\mid R_u(o)\right]\to 0,
	\qquad\text{as }u\to\infty.
	\]
\end{theorem}

\begin{proof}
	For fixed $k$, let every peak in a size-$k$ cluster send one unit of mass
	to the root of that cluster.  The mass-transport principle gives
	\begin{equation}\label{eq:root-peak-mass-transport}
		\begin{split}
		w_k^{\rm peak}
		&= \E\left[
		M_u(o)\mathbbm{1}\{|C_u(o)|=k\}\mathbbm{1}\{R_u(o)\}
		\right]\\
		&=
		\rho_u
		\E\left[
		M_u(o)\mathbbm{1}\{|C_u(o)|=k\}\mid R_u(o)
		\right].
	\end{split}
	\end{equation}
	Summing over $k$ yields
	\begin{equation}\label{eq:peak-root-intensity-relation}
		\rho_u^{\rm peak}
		=
		\rho_u\E\left[M_u(o)\mid R_u(o)\right].
	\end{equation}
	
	Let
	\[
	\mu_u=\E\left[M_u(o)\mid R_u(o)\right],\quad
	a_{k,u}
	=
	\E\left[
	(M_u(o)-1)\mathbbm{1}\{|C_u(o)|=k\}\mid R_u(o)
	\right].
	\]
	Then $a_{k,u}\geq0$ and
	\[
	\sum_{k\geq1}a_{k,u}=\mu_u-1.
	\]
	Taking the ratio between \eqref{eq:root-peak-mass-transport} and \eqref{eq:peak-root-intensity-relation} yields
	\[
	\P(S_u^{\rm peak}=k)
	=
	\frac{\E\left[\mathbbm{1}\{|C_u(o)|=k\}\mid R_u(o)
		\right]+a_{k,u}}{\mu_u} = \frac{\P(S_u=k) + a_{k,u}}{\mu_u}.
	\]
	Therefore,
	\begin{align*}
		&\quad \sum_{k=1}^{\infty} \left|
		\P(S_u^{\rm peak}=k)
		-\P(S_u=k)
		\right|\\
		&\leq
		\frac{1}{\mu_u}
		\left\{
		\sum_{k=1}^\infty a_{k,u}
		+
		(\mu_u-1)\sum_{k=1}^\infty \P(S_u=k)
		\right\}=
		\frac{\mu_u-1}{\mu_u}.
	\end{align*}
	Equation~\eqref{eq:peak-root-intensity-relation} gives the final equality
	in \eqref{eq:root-peak-TV-bound}.
\end{proof}

The following result shows that the total variation distance between the exact and peak-based cluster size distributions of Gaussian random fields vanishes exponentially fast as $u\to \infty$.
\begin{corollary}[Gaussian high-threshold equivalence]
	\label{cor:gaussian-root-peak}
	Let $\{X_t:t\in\Z^d\}$ be a stationary, centered, unit-variance Gaussian
	field. Suppose that 
	Assumption~\ref{ass:finite-cluster} holds for all sufficiently large $u$,
	and that
	\begin{equation}\label{eq:r*}
	r_*:=
	\max_{s\in\mathcal N(o)}
	\operatorname{Corr}(X_o,X_s)\in (-1, 1).
    \end{equation}
	Then there exists $\alpha>0$ such that
	\[
	\sum_{k=1}^{\infty} \left|
	\P(S_u^{\rm peak}=k)
	-\P(S_u=k)
	\right| \le e^{-\alpha u^2}
	\to 0,
	\qquad \text{\rm as } u\to\infty.
	\]
\end{corollary}

\begin{proof}
	Write $\Psi(x)=1-\Phi(x)$ for the standard Gaussian tail function. Let
	\[
	q_u=\P(X_o>u)=\Psi(u)
	\]
	and
	\[
	b_u
	=
	\P\left(
	X_o>u,\
	\max_{s\in\mathcal N(o)}X_s>u
	\right).
	\]
	An exceedance at the origin with no neighboring exceedance is a singleton
	cluster. The origin is then both its root and a local maximum.
	Therefore,
	\[
	q_u-b_u\leq\rho_u.
	\]
	Moreover, Theorem~\ref{thm:root-peak-comparison} gives
	$\rho_u\leq\rho_u^{\rm peak}$, while
	$\rho_u^{\rm peak}\leq q_u$. Hence
	\[
	q_u-b_u
	\leq\rho_u
	\leq\rho_u^{\rm peak}
	\leq q_u.
	\]
	Using Theorem~\ref{thm:root-peak-comparison} again,
	\begin{equation}\label{eq:gaussian-TV-short-bound}
		\sum_{k=1}^{\infty} \left|
		\P(S_u^{\rm peak}=k)
		-\P(S_u=k)
		\right|
		\leq
		1-\frac{\rho_u}{\rho_u^{\rm peak}}
		\leq
		\frac{b_u}{q_u}.
	\end{equation}
	
	For every
	$s\in\mathcal N(o)$,
	\[
	\P(X_o>u,X_s>u)
	\leq
	\Psi\left(
	\sqrt{2/(1+r_*)}\,u
	\right).
	\]
	Thus, by the union bound,
	\[
	\frac{b_u}{q_u}
	\leq
	|\mathcal N(o)|
	\frac{
		\Psi\left(
		\sqrt{2/(1+r_*)}\,u
		\right)}
	{\Psi(u)}
	\to 0,
	\]
	where the convergence follows from Mills' ratio because $r_*<1$.
	The conclusion now follows from
	\eqref{eq:gaussian-TV-short-bound} by choosing $0<\alpha<(1-r_*)/[2(1+r_*)]$.
\end{proof}

\section{Monte Carlo Estimation}
\label{sec:MC}
Using independent field realizations, we estimate both exact and
peak-based cluster size distributions by spatial Monte Carlo
averaging.  A deterministic guard band removes finite-window bias for all
cluster sizes up to a cutoff $K$.  Throughout this section the field is
stationary, Assumption~\ref{ass:finite-cluster} holds, and the relevant
normalizers $\rho_u$ and $\rho_u^{\rm peak}$ are positive.

For a fixed cutoff $K\geq1$, define
\begin{align*}
	\rho_K:=\sum_{k=1}^K w_k,
	\qquad \tau_K:=\sum_{k>K}w_k=\rho_u-\rho_K.
\end{align*}
By \eqref{eq:peak-event} and \eqref{eq:peak-intensity-local}, the peak normalizer $\rho_u^{\rm peak}$ is a finite-dimensional probability
involving only a site and its neighbors. It is therefore assumed to be evaluated directly from their joint distribution. The theoretical results below treat this value as exact.

\subsection{Guarded spatial estimation}

\subsubsection{Exact computation with guard bands}
\label{sec:mc-guard}

Let $G$ be the lattice graph and $d_G$ its graph distance.  For
$A\subset\Z^d$ and $r\geq0$, define 
\begin{equation*}
	A^{\oplus r}:=\left\{s\in\Z^d: \inf_{a\in A}d_G(s,a)\leq r\right\},
\end{equation*}
and write it simply $t^{\oplus r}$ if $A=\{t\}$. Thus $A^{\oplus r}\setminus A$ is the distance-$r$ guard band around
$A$.

For a finite simulation window $V\subset\Z^d$ and $t\in V$, let $C_u(t|V)$ be the connected component of $t$ restricted on $E_u\cap V$, and set $C_u(t|V)=\varnothing$ when $X_t\leq u$.
Equivalently,
\begin{equation*}
	C_u(t|V)
	=
	\left\{s\in E_u\cap V:
		\text{there is an excursion path in $V$ connects $t$ to $s$}
	\right\}.
\end{equation*}
The component $C_u(t|V)$ uses only observations in $V$ and may be smaller
than the full-lattice component $C_u(t)$ if the latter is not closed inside
$V$.

\begin{lemma}[Exact finite determination]
	\label{lem:mc-guard-band}
	If $t^{\oplus K}\subset V$, then, for every $1\leq k\leq K$,
	\begin{align*}
		\{|C_u(t)|=k\}=\{|C_u(t|V)|=k\},
		\qquad
		\{|C_u(t)|>K\}=\{|C_u(t|V)|>K\}.
	\end{align*}
	Consequently, observations on $V=\Lambda^{\oplus K}$ determine exactly,
	for every $t\in\Lambda$, whether $t$ is a nonexceedance, belongs to a
	cluster of size $1,\ldots,K$, or belongs to a cluster of size greater than
	$K$.  They also determine $A_u^{\rm peak}(t)$.
\end{lemma}

\begin{proof}
	If $|C_u(t)|=k\leq K$, every vertex of the component lies in
	$t^{\oplus(k-1)}$, and every exterior neighbor lies in
	$t^{\oplus k}\subset V$.  Hence the component closes inside $V$ and
	$C_u(t\mid V)=C_u(t)$.  If $|C_u(t)|>K$, a breadth-first search finds
	$K+1$ connected exceedance sites in $t^{\oplus K}\subset V$, so
	$|C_u(t\mid V)|>K$. The converse follows from
	$C_u(t\mid V)\subset C_u(t)$.
	Finally, $A_u^{\rm peak}(t)$ depends only on $X_t$ and the values at
	$\mathcal N(t)\subset t^{\oplus 1}\subset t^{\oplus K}$.
\end{proof}

Thus $\Lambda$ is the scoring window and $V=\Lambda^{\oplus K}$ is the larger simulation window that provides the required guard band. Connected components are labeled with the usual lattice adjacency on $V$, without periodic wrap-around.

\subsubsection{Single-origin estimators}
\label{sec:mc-estimators}

For $1\leq k\leq K$, define the exact and peak site scores
\begin{equation}\label{eq:xi}
	\xi_{t,k}
	:=\frac1k\mathbbm{1}\{|C_u(t)|=k\},
	\qquad
	\xi_{t,k}^{\rm peak}
	:=\mathbbm{1}\{A_u^{\rm peak}(t),\ |C_u(t)|=k\}.
\end{equation}
By stationarity, \eqref{eq:w-inside-w} and \eqref{eq:w-peak} give
\begin{equation*}
	\E\xi_{t,k}=w_k,
	\qquad
	\E\xi_{t,k}^{\rm peak}=w_k^{\rm peak}.
\end{equation*}
The factor $1/k$ removes the size bias caused by observing a size-$k$
cluster from each of its $k$ sites.  No such factor is used for peaks, since
each peak is a sampling location.

Let $X^{(1)},\ldots,X^{(M)}$ be independent field realizations. A superscript $(m)$ denotes the quantity computed
from realization $m$. The single-origin estimators for $w_k$ and $w_k^{\rm peak}$, obtained by simulating on $V_o=o^{\oplus K}$, are respectively
\begin{equation*}
	\begin{aligned}
		\widehat w_k^{\rm SO}
		&=\frac1{Mk}\sum_{m=1}^M
		\mathbbm{1}\{|C_u^{(m)}(o\mid V_o)|=k\},\\
		\widehat w_k^{{\rm peak},{\rm SO}}
		&=\frac1M\sum_{m=1}^M
		\mathbbm{1}\{A_u^{{\rm peak},(m)}(o),\
		|C_u^{(m)}(o\mid V_o)|=k\}.
	\end{aligned}
\end{equation*}


\subsubsection{Spatial estimators}

The single-origin estimator uses only the excursion component associated with one prespecified site and therefore leaves most of the information in a simulated field unused. We next introduce a spatial procedure that extracts more information from each realization by using all sites of a finite scoring window $\Lambda\subset\Z^d$. To ensure that cluster sizes up to $K$ are determined exactly, the field is simulated on the enlarged window $V=\Lambda^{\oplus K}$. For one realization define
\begin{align}
	W_{m,k}
	&=\frac{1}{k|\Lambda|}\sum_{t\in\Lambda}
	\mathbbm{1}\{|C_u^{(m)}(t|V)|=k\},
	\label{eq:MC2}\\
	W_{m,k}^{{\rm peak}}
	&=\frac{1}{|\Lambda|}\sum_{t\in\Lambda}
	\mathbbm{1}\{A_u^{{\rm peak},(m)}(t),\ |C_u^{(m)}(t|V)|=k\}.
	\label{eq:MC3}
\end{align}
Their averages over independent realizations give the following \textit{spatial Monte Carlo estimators} for $w_k$ and $w_k^{\rm peak}$:
\begin{equation}\label{eq:wk-estimate}
	\widehat w_k:=\frac1M\sum_{m=1}^M W_{m,k},
	\qquad
	\widehat w_k^{\rm peak}:=\frac1M\sum_{m=1}^M W_{m,k}^{\rm peak}.
\end{equation}
The spatial peak-based probability estimator is therefore
\begin{equation*}
	\widehat p_k^{\rm peak}
	:=\frac{\widehat w_k^{\rm peak}}{\rho_u^{\rm peak}}.
\end{equation*}
Because the full normalizer $\rho_u^{\rm peak}$ is used, the estimates reported for $k\leq K$ are not renormalized and need not sum to one. The choice $\Lambda=\{o\}$ recovers the single-origin estimator; larger
windows use more of each simulated realization.

Rather than determining $C_u(t| V)$ separately for every
$t\in\Lambda$, we label the connected components of the excursion set
$E_u\cap V$ once. The resulting component labels and sizes provide all the quantities needed to evaluate the sums in \eqref{eq:MC2}--\eqref{eq:MC3}. Let
$\mathscr C_u^{(m)}(V)$ denote the collection of labeled excursion components in
realization $m$.  For $D\in\mathscr C_u^{(m)}(V)$, let
\[
a_D=|D\cap\Lambda|,
\qquad
b_D=\#\{t\in D\cap\Lambda:
A_u^{{\rm peak},(m)}(t)\}.
\]
Only components satisfying $D\cap\Lambda\neq\varnothing$ need to be processed. Lemma~\ref{lem:mc-guard-band} guarantees that every component intersecting
$\Lambda$ is classified correctly as having size $k\leq K$ or size greater
than $K$, and hence gives the following algorithm.
\begin{enumerate}[label=(\roman*)]
	\item Generate the field on $V=\Lambda^{\oplus K}$, and label its connected excursion components on $E_u\cap V$.
	\item For each component $D$ with size $k\leq K$, add $a_D/k$ to the exact size-$k$ total and add $b_D$ to the peak size-$k$ total.
	\item For each component $D$ with size greater than $K$, add $a_D$ to the exact overflow total.
	\item Divide all totals by $|\Lambda|$ and average them over the $M$
	independent realizations.
\end{enumerate}
Applying this algorithm, the scores in \eqref{eq:MC2}--\eqref{eq:MC3} become
\begin{align}\label{eq:W2}
	W_{m,k}
	=
	\frac1{|\Lambda|}
	\sum_{\substack{D\in\mathscr C_u^{(m)}(V)\\
			|D|=k}}
	\frac{|D\cap\Lambda|}{k}, \qquad 
	W_{m,k}^{\rm peak}
	=
	\frac1{|\Lambda|}
	\sum_{\substack{D\in\mathscr C_u^{(m)}(V)\\
			|D|=k}}
	b_D.
\end{align}
In general, a size-$k$ component $D$ contributes $|D\cap\Lambda|/|D|$ to the unnormalized exact total. For exact law tail control, also define and record
\begin{equation}\label{eq:mc-overflow-estimators}
	\begin{split}
	T_m&:=\frac1{|\Lambda|}\sum_{t\in\Lambda}
	\mathbbm{1}\{|C_u^{(m)}(t\mid V)|>K\} = \frac1{|\Lambda|}
	\sum_{\substack{D\in\mathscr C_u^{(m)}(V)\\|D|>K}}
	|D\cap\Lambda|,\\
	\widehat r_K&:=\frac1M\sum_{m=1}^M T_m, \qquad r_K:=\P(|C_u(o)|>K).
\end{split}
\end{equation}


\begin{remark}[Peak-based estimation]
	\label{remark:MC-peak}
	The estimator \eqref{eq:MC3} is an unconditional spatial average under
	the original field law, not simulation conditional on a peak.  For a
	nonstationary field, the single-origin construction remains valid at a
	fixed site $t$, with the quantities
	$w_k^{\rm peak}(t)$ and $\rho_u^{\rm peak}(t)$ depending on $t$.
\end{remark}

\subsection{Unbiasedness, consistency, and truncation control}
\label{sec:mc-unbiased}

Define 
\begin{equation*}
	\widehat\rho_K:=\sum_{k=1}^K\widehat w_k.
\end{equation*}

\begin{theorem}[Unbiasedness and consistency]
	\label{thm:mc-unbiased}
	Suppose that the field is stationary and that
	Assumption~\ref{ass:finite-cluster} holds.  For every finite nonempty
	$\Lambda$, every $M\geq1$, and every $1\leq k\leq K$,
	\begin{equation}\label{eq:mc-unbiasedness}
		\E\widehat w_k=w_k,
		\qquad
		\E\widehat w_k^{\rm peak}=w_k^{\rm peak},
		\qquad
		\E\widehat\rho_K=\rho_K,
		\qquad
		\E\widehat r_K=r_K.
	\end{equation}
	All four estimators converge almost surely to their targets as
	$M\to\infty$.  If $\rho_K>0$ and $\rho_u^{\rm peak}>0$, then
	\begin{equation}\label{eq:mc-ratio-consistency}
		\frac{\widehat w_k}{\widehat\rho_K}
		\to\frac{w_k}{\rho_K},
		\qquad
		\widehat p_k^{\rm peak}
		\to
		\P(S_u^{\rm peak}=k), \quad {\rm a.s}.
	\end{equation}
	Moreover, because $\rho_u^{\rm peak}$ is nonrandom,
	$\widehat p_k^{\rm peak}$ is unbiased for
	$\P(S_u^{\rm peak}=k)$ for every $M$.
\end{theorem}

\begin{proof}
	By Lemma~\ref{lem:mc-guard-band}, the cluster-size and overflow
	classifications computed on $V=\Lambda^{\oplus K}$ agree with their
	full-lattice counterparts at every site of $\Lambda$.  Therefore,
	stationarity and linearity of expectation give
	\begin{align*}
		\E W_{m,k}
		&=\frac1{|\Lambda|}\sum_{t\in\Lambda}\E\xi_{t,k}
		=w_k, \quad \E W_{m,k}^{\rm peak}
		=\frac1{|\Lambda|}\sum_{t\in\Lambda}
		\E\xi_{t,k}^{\rm peak}
		=w_k^{\rm peak}.\\
		\E T_m
		&=\frac1{|\Lambda|}\sum_{t\in\Lambda}
		\P(C_u(t)>K) = 
		r_K.
	\end{align*}
	Averaging over $m$ proves \eqref{eq:mc-unbiasedness}.  
	
	For each estimator, the replication-level quantities are independent and
	identically distributed across $m$ and are bounded.  The strong law of
	large numbers proves the consistency in \eqref{eq:mc-ratio-consistency}. Finally, $\E\widehat p_k^{\rm peak}=w_k^{\rm peak}/\rho_u^{\rm peak}$ follows
	from the deterministic denominator.
\end{proof}

The exact ratio in \eqref{eq:mc-ratio-consistency} estimates the truncated
law $w_k/\rho_K$, rather than $\P(S_u=k)=w_k/\rho_u$. 	Using $\P(|C_u(o)|=j)=j w_j$, we have
\[
r_K
=\sum_{j>K}\P(|C_u(o)|=j)
=\sum_{j>K}j w_j
\geq(K+1)\sum_{j>K}w_j=(K+1)\tau_K.
\]
Thus, we obtain the following result that
controls the difference using the observable overflow probability $r_K$:
	\begin{align}\label{eq:tau-B}
		0\leq\tau_K
		\leq B_K:=\frac{r_K}{K+1}.
	\end{align}
For many fields, including Gaussian fields under suitable regularity conditions, large excursion clusters are rare, and the tail intensity $r_K$ and hence $\tau_K$ decreases rapidly to zero as $K\to\infty$; see Remark \ref{remark:ld-mc-truncation}.

The effect of the exact tail on normalization can be expressed as
\begin{equation*}
	\rho_K\leq\rho_u\leq\rho_K+B_K,
	\qquad
	\frac{w_k}{\rho_K+B_K}
	\leq\P(S_u=k)
	\leq\frac{w_k}{\rho_K},
	\quad 1\leq k\leq K.
\end{equation*}
To describe the complete truncation error, suppose $\rho_K>0$, and define the renormalized truncated mass functions,
extended by zero for $k>K$, by
\[
p_{k,K}=
\begin{cases}
	w_k/\rho_K,&k\leq K,\\
	0,&k>K.
\end{cases}
\]
Then we have
\begin{equation}\label{eq:mc-exact-TV-truncation}
	\begin{split}
	\sum_{k=1}^{\infty} \left|
	p_{k,K}
	-\P(S_u=k)
	\right| &= \sum_{k=1}^K \left(\frac{w_k}{\rho_K} - \frac{w_k}{\rho_u}\right)+ \sum_{k>K} \frac{w_k}{\rho_u}\\
	&= \left(1-\frac{\rho_K}{\rho_u}\right) + \frac{\tau_K}{\rho_u} =\frac{2\tau_K}{\rho_u}
	\leq\frac{2B_K}{\rho_K+B_K}.
	\end{split}
\end{equation}

\begin{corollary}[Consistency for the exact distribution]
	\label{cor:mc-two-stage-consistency}
	Whenever $\widehat\rho_K>0$, define
	\[
	\widehat p_{k,K,M}
	=
	\begin{cases}
		\widehat w_k/\widehat\rho_K,&1\leq k\leq K,\\
		0,&k>K.
	\end{cases}
	\]
	Under the assumptions of Theorem~\ref{thm:mc-unbiased}, assume also that
	$\rho_K>0$.  Then
	\begin{equation*}
		\begin{split}
		\lim_{M\to\infty}\sum_{k=1}^{\infty} \left|
		\widehat p_{k,K,M}
		-\P(S_u=k)
		\right|
		&=\frac{2\tau_K}{\rho_u}, \quad \text{\rm a.s.},\\
		\lim_{K\to\infty}\lim_{M\to\infty}
		\sum_{k=1}^{\infty} \left|
		\widehat p_{k,K,M}
		-\P(S_u=k)
		\right|&=0, \quad \text{\rm a.s.}
		\end{split}
	\end{equation*}
\end{corollary}

\begin{proof}
	For fixed $K$, Theorem~\ref{thm:mc-unbiased} gives simultaneous
	almost sure convergence of the finitely many estimated weights and their
	sum.  Hence 
	\[
	\lim_{M\to\infty}\sum_{k=1}^{\infty} \left|
	\widehat p_{k,K,M}
	- p_{k,K}
	\right| =0, \qquad \text{a.s.}
	\]
	Eq.~\eqref{eq:mc-exact-TV-truncation} yields the first limit.
	Finally, $\tau_K\downarrow0$ because
	$\rho_u=\sum_{k\geq1}w_k<\infty$, proving the second limit.
\end{proof}


\begin{comment}
If the upper bound at $K$ is too conservative, choose $L>K$ as a new guard
band.  Then
\begin{equation*}
	\sum_{j=K+1}^{L}w_j
	\leq\tau_K
	\leq\sum_{j=K+1}^{L}w_j+\frac{r_L}{L+1}.
\end{equation*}
The finite sum is estimated unbiasedly by
$\sum_{j=K+1}^{L}\widehat w_j$, while $r_L$ is estimated by the spatial
overflow score at cutoff $L$.  Increasing $L$ therefore replaces part of the
upper allowance by explicitly estimated cluster weights.
\end{comment}

\subsection{Monte Carlo variance of the spatial estimators}
\label{sec:mc-spatial-variance}

We now quantify the Monte Carlo error of the spatial estimator $\widehat w_k$ and make its dependence on the size and geometry of $\Lambda$ explicit. The same
argument applies to $\widehat w_k^{\rm peak}$ after replacing
$\xi_{t,k}$ by $\xi_{t,k}^{\rm peak}$. 

For $h\in\Z^d$, define 
\[
\gamma_k(h)
=\operatorname{Cov}(\xi_{o,k}, \xi_{o+h,k}),
\]
and 
\begin{equation*}
	N_\Lambda(h)
	=\#\{(s,t)\in\Lambda^2:t-s=h\}
	=|\Lambda\cap(\Lambda-h)|,
	\qquad
	a_\Lambda(h)=\frac{N_\Lambda(h)}{|\Lambda|}.
\end{equation*}

\begin{proposition}[Window-dependent spatial variance]
	\label{prop:mc-window-variance}
	Suppose that the field is stationary.  For every finite nonempty
	$\Lambda$, every $1\leq k\leq K$, 
	\begin{equation}\label{eq:mc-window-variance}
		\begin{split}
		\operatorname{Var}(\widehat w_k)
		=\frac1{M|\Lambda|^2}
		\sum_{h\in\Z^d}N_\Lambda(h)\gamma_k(h)
		=\frac1{M|\Lambda|}
		\sum_{h\in\Z^d}a_\Lambda(h)\gamma_k(h).
	\end{split}
	\end{equation}
	If $\sum_{h\in\Z^d}|\gamma_k(h)|<\infty$,
	then
	\begin{equation}\label{eq:mc-window-variance-bound}
		\operatorname{Var}(\widehat w_k)
		\leq\frac1{M|\Lambda|}
		\sum_{h\in\Z^d}|\gamma_k(h)|.
	\end{equation}
	Consequently, the variance is $O((M|\Lambda|)^{-1})$ for fixed $k$.
\end{proposition}

\begin{proof}
	Lemma~\ref{lem:mc-guard-band} allows the guarded scores computed on $V$
	to be replaced by their full-lattice counterparts. By \eqref{eq:xi} and \eqref{eq:MC2}, 
	\[
	W_{m,k}
	=\frac1{|\Lambda|}\sum_{t\in\Lambda}\xi_{t,k}^{(m)}.
	\]
	Stationarity gives
	\begin{align*}
		\operatorname{Var}(W_{m,k})
		=\frac1{|\Lambda|^2}
		\sum_{s,t\in\Lambda}
		\operatorname{Cov}(\xi_{s,k},\xi_{t,k})=\frac1{|\Lambda|^2}
		\sum_{h\in\Z^d}N_\Lambda(h)\gamma_k(h).
	\end{align*}
	Averaging $M$ independent realizations divides this variance by $M$,
	proving \eqref{eq:mc-window-variance}.  The bound in \eqref{eq:mc-window-variance-bound} follows from $0\leq a_\Lambda(h)\leq1$ and the triangle inequality.
\end{proof}

For a rectangular scoring window $\Lambda=\prod_{i=1}^d\{1,\ldots,n_i\}$ with $|\Lambda|=\prod_{i=1}^dn_i$, the overlap count is $N_\Lambda(h)=\prod_{i=1}^d(n_i-|h_i|)_+$ where $(x)_+:=\max\{x,0\}$.
Along regular expanding windows $\{\Lambda_n\}$ satisfying
$|\Lambda_n|\to\infty$ and $a_{\Lambda_n}(h)\to 1$ for every fixed $h$,
we have
\begin{equation*}
	M|\Lambda_n|\operatorname{Var}(\widehat w_{k,n})
	\to \sum_{h\in\Z^d}\gamma_k(h),
	\qquad \text{\rm as } n\to\infty.
\end{equation*}

\section{Cluster Tails and High-Threshold Large Deviations}
\label{sec:large-deviations}

We consider two complementary regimes: $k\to\infty$ at fixed threshold and
$u\to\infty$ at fixed cluster size.  The first yields conservative
exponential tail bounds useful for Monte Carlo truncation; the second yields
exact $u^2$-scale logarithmic asymptotics for Gaussian fields.

One comparison will be used in both regimes.  A size-$k$ cluster contributes once to the exact weight and once for each of its weak peaks to the peak weight.
Every finite size-$k$ cluster has between one and $k$ weak peaks, so
\begin{equation}\label{eq:ld-weight-sandwich}
	w_k\leq w_k^{\rm peak}\leq k w_k.
\end{equation}
In the fixed-threshold regime, $(\log k)/k\to0$; in the high-threshold
regime, $k$ is fixed and $(\log k)/u^2\to0$.  Thus the factor is negligible
on the relevant exponential scale in either regime.  

\subsection{Large-cluster bounds at a fixed threshold}
\label{sec:ld-large-k}

In this subsection, $u$ is fixed and the field is stationary.  We give a
sufficient condition for exponential cluster tails and then verify it for a
class of Gaussian fields.

\subsubsection{A general exponential bound}
\label{sec:ld-general-bound}

Recall that $\mathcal C_k^{\rm inside}$ is the collection of connected
sets of $k$ sites that contain the origin; and $|\mathcal C_k^{\rm inside}|$ denote the number of clusters it contains. Define the maximal vertex degree of the lattice graph by
\[
\Delta:=|\mathcal N(o)|.
\]
A depth-first contour traversal of a deterministic spanning tree of
$D\in\mathcal C_k^{\rm inside}$ has length $2(k-1)$ and determines $D$.
Since at each step
there are at most $\Delta$ choices,
\begin{equation}\label{eq:ld-animal-count}
|\mathcal C_k^{\rm inside}|
\leq \Delta^{2(k-1)}\leq\beta^k,
\qquad \text{\rm where } \beta:=\Delta^2.
\end{equation}
For nearest-neighbor adjacency on $\Z^d$, $\Delta=2d$ and we have
$\beta=4d^2$.  In $d=1$, the exact count is
$|\mathcal C_k^{\rm inside}|=k$, so the sharper choice $\beta=2$ is valid.

\begin{assumption}[Uniform cluster exceedance bound]
	\label{ass:ld-connected-exceedance}
	There exists $\psi_u>0$ such that, for every $k\geq1$,
	\begin{equation}\label{eq:ld-joint-exceedance}
		\sup_{D\in\mathcal C_k^{\rm inside}}
		\P(X_t>u\text{ for every }t\in D)
		\leq e^{-\psi_u k}.
	\end{equation}
\end{assumption}


The next theorem expresses the competition between the exponential number
of possible cluster shapes and the probability of simultaneous exceedances
on any one shape.

\begin{theorem}[Exponential bound for large clusters]
	\label{thm:ld-general}
	Suppose Assumption~\ref{ass:ld-connected-exceedance} holds and
	$\psi_u>\log \beta$.  Set
	\begin{equation}\label{eq:ld-cu}
		c_u:=\psi_u-\log \beta>0.
	\end{equation}
	Then, for every $k\geq1$,
	\begin{align}
		\P\{|C_u(o)|\geq k\}
		&\leq e^{-c_u k}.
		\label{eq:ld-site-tail}
	\end{align}
	In particular, Assumption~\ref{ass:finite-cluster} holds.  If
	$\rho_u,\rho_u^{\rm peak}>0$, then
		\begin{align}
		\P(S_u\geq k)
		&\leq\frac{e^{-c_u k}}{k\rho_u},
		\label{eq:ld-exact-tail}\\
		\P(S_u^{\rm peak}\geq k)
		&\leq\frac{e^{-c_u k}}{\rho_u^{\rm peak}}.
		\label{eq:ld-peak-tail}
	\end{align}
\end{theorem}

\begin{proof}
	If $|C_u(o)|\geq k$, a breadth-first exploration of $C_u(o)$ produces a
	connected subset $D\in\mathcal C_k^{\rm inside}$ on which every site
	exceeds $u$.  Hence, by \eqref{eq:ld-animal-count},
	Assumption~\ref{ass:ld-connected-exceedance}, and the union bound,
	\[
	\P\{|C_u(o)|\geq k\}
	\leq
	\sum_{D\in\mathcal C_k^{\rm inside}}
	\P(X_t>u\text{ for every }t\in D)
	\leq \beta^k e^{-\psi_u k}
	=e^{-c_u k}.
	\]
	Letting $k\to\infty$ shows that
	$\P\{|C_u(o)|=\infty\}=0$.  By stationarity and the countability of
	$\Z^d$, every excursion cluster is therefore finite almost surely under
	the assumptions of the theorem.
	Furthermore, \eqref{eq:w-inside-w} yields
	\[
	\P(S_u\geq k)
	=\frac1{\rho_u}\sum_{j\geq k}w_j
	\leq\frac1{k\rho_u}\sum_{j\geq k}jw_j
	=\frac{\P\{|C_u(o)|\geq k\}}{k\rho_u}.
	\]
	For the peak-based law,
	\[
	\P(S_u^{\rm peak}\geq k)
	=
	\frac{\P\{A_u^{\rm peak}(o),\ |C_u(o)|\geq k\}}
	{\rho_u^{\rm peak}}
	\leq
	\frac{\P\{|C_u(o)|\geq k\}}{\rho_u^{\rm peak}}.
	\]
	Combining these inequalities with \eqref{eq:ld-site-tail} proves
	\eqref{eq:ld-exact-tail} and \eqref{eq:ld-peak-tail}.
\end{proof}

The condition $\psi_u>\log \beta$ has a direct interpretation:
$\psi_u$ is the probabilistic cost per site of maintaining simultaneous
exceedances, whereas $\log \beta$ is an upper bound on the combinatorial growth
per site allowed by the count of connected shapes.  Large clusters have an
exponentially small probability whenever the former dominates the latter.

\subsubsection{Gaussian truncation bounds}
\label{sec:ld-gaussian}

For a Gaussian field, the joint exceedance condition has a
simple sufficient criterion.

\begin{corollary}[Gaussian large cluster bound]
	\label{cor:ld-gaussian}
	Let $X$ be a centered stationary Gaussian field with covariance
	$\Gamma(h):=\operatorname{Cov}(X_o,X_h)$, and suppose
	\begin{equation*}
		0<\kappa:=\sum_{h\in\Z^d}|\Gamma(h)|<\infty.
	\end{equation*}
	If $u>\sqrt{2\kappa\log \beta}$,
	then \eqref{eq:ld-site-tail}--\eqref{eq:ld-peak-tail} hold with
	\begin{equation*}
		c_u=\frac{u^2}{2\kappa}-\log \beta>0.
	\end{equation*}
\end{corollary}

\begin{proof}
	Fix $D\in\mathcal C_k^{\rm inside}$ and set
	$Y_D=\sum_{t\in D}X_t$.  By stationarity,
	\[
	\operatorname{Var}(Y_D)
	=\sum_{s,t\in D}\Gamma(t-s)
	\leq
	\sum_{s\in D}\sum_{h\in\Z^d}|\Gamma(h)|
	=\kappa k.
	\]
	The event $\{X_t>u\text{ for every }t\in D\}$ implies $Y_D>ku$.
	If $\operatorname{Var}(Y_D)=0$, this event has probability zero.
	Otherwise, the Gaussian Chernoff bound gives
	\[
	\P(X_t>u\text{ for every }t\in D)
	\leq\P(Y_D>ku)
	\leq
	\exp\!\left\{-\frac{k^2u^2}
	{2\operatorname{Var}(Y_D)}\right\}
	\leq\exp\!\left(-\frac{ku^2}{2\kappa}\right).
	\]
	Thus Assumption~\ref{ass:ld-connected-exceedance} holds with
	$\psi_u=u^2/(2\kappa)$, and
	Theorem~\ref{thm:ld-general} applies.
\end{proof}

\begin{remark}[Implication for Monte Carlo truncation]
	\label{remark:ld-mc-truncation}
	Under Theorem~\ref{thm:ld-general}, taking $k=K+1$ gives
	\begin{equation*}
		r_K=\P\{|C_u(o)|>K\}
		\leq e^{-c_u(K+1)}.
	\end{equation*}
	Consequently, by \eqref{eq:tau-B},
	\begin{equation*}
		\tau_K\leq B_K
		\leq\frac{e^{-c_u(K+1)}}{K+1},
	\end{equation*}
	and the total variation error in
	\eqref{eq:mc-exact-TV-truncation} obeys
	\begin{equation*}
		\sum_{k=1}^{\infty} \left|
		p_{k,K}
		-\P(S_u=k)
		\right|
		\leq
		\frac{2e^{-c_u(K+1)}}{(K+1)\rho_u}.
	\end{equation*}
	The peak-based probability is bounded directly by
	\begin{equation*}
		\P(S_u^{\rm peak}>K)
		\leq\frac{e^{-c_u(K+1)}}{\rho_u^{\rm peak}}.
	\end{equation*}
	Thus the bound provides a conservative guide for selecting
	the Monte Carlo cutoff $K$.
\end{remark}

\subsection{Gaussian high-threshold large deviations at fixed size}
\label{sec:ld-high-threshold}

We now fix $k$ and let $u\to\infty$.  In this subsection, let $X$ be a
centered stationary Gaussian field with variance $\sigma^2$, and assume that its covariance matrix on every finite set of distinct sites
is positive definite.  We also assume that
Assumption~\ref{ass:finite-cluster} holds for all sufficiently large $u$.  To display the dependence on the threshold,
write the weights as $w_k(u)$ and $w_k^{\rm peak}(u)$ in this subsection.

For $D\in\mathcal C_k^{\rm root}$, note that $D^{\oplus 1}=D\cup\mathcal N(D)$; let $\Sigma_{D^{\oplus 1}}$ be the covariance matrix of
$(X_t)_{t\in D^{\oplus 1}}$, and define the closed polyhedron
\begin{equation}\label{eq:ld-shape-polyhedron}
	\mathcal K_D
	:=
	\left\{
	x=(x_t)_{t\in D^{\oplus 1}}:
	x_t\geq1\ \text{for }t\in D,
	\,
	x_s\leq1\ \text{for }s\in\mathcal N(D)
	\right\}.
\end{equation}
The rate associated with the rooted shape $D$ is
\begin{equation*}
	I_D
	:=
	\inf_{x\in\mathcal K_D}
	\frac12x^\top\Sigma_{D^{\oplus 1}}^{-1}x,
\end{equation*}
and the size-$k$ rate is
\begin{equation*}
	I_k:=\min_{D\in\mathcal C_k^{\rm root}}I_D.
\end{equation*}
The minimum is over a finite collection.  Thus $I_k$ is the Gaussian cost
of the least unlikely rooted cluster shape of size $k$.  The exterior
constraints in \eqref{eq:ld-shape-polyhedron} certify that $D$ is the complete
cluster, rather than merely a connected set of exceedance sites, and in
general they cannot be omitted from the optimization.  

\begin{theorem}[Gaussian high-threshold large deviations]
	\label{thm:ld-high-threshold}
	Under the preceding assumptions, for every fixed $k\geq1$,
	\begin{equation}\label{eq:ld-weight-high-u}
		\lim_{u\to\infty}
		-\frac1{u^2}\log w_k(u)
		=
		\lim_{u\to\infty}
		-\frac1{u^2}\log w_k^{\rm peak}(u)
		=I_k.
	\end{equation}
	Moreover,
	\begin{equation}\label{eq:ld-normalizer-high-u}
		\lim_{u\to\infty}
		-\frac1{u^2}\log\rho_u
		=
		\lim_{u\to\infty}
		-\frac1{u^2}\log\rho_u^{\rm peak}
		=\frac1{2\sigma^2}.
	\end{equation}
	Consequently,
	\begin{equation}\label{eq:ld-pmf-high-u}
		\begin{split}
			\lim_{u\to\infty}
			-\frac1{u^2}\log\P(S_u=k)
			&=
			\lim_{u\to\infty}
			-\frac1{u^2}\log\P(S_u^{\rm peak}=k)=I_k-\frac1{2\sigma^2}.
		\end{split}
	\end{equation}
\end{theorem}

\begin{proof}
	Fix $D\in\mathcal C_k^{\rm root}$.  The vectors
	$u^{-1}(X_t)_{t\in D^{\oplus 1}}$ satisfy a finite-dimensional Gaussian
	large deviation principle with speed $u^2$ and rate function
	\[
	Q_D(x)=\frac12x^\top\Sigma_{D^{\oplus 1}}^{-1}x.
	\]
	By \eqref{eq:C_u(o)}, the event that $D$ is the excursion cluster rooted at the origin is $\{C_u(o)=D\}$. Let
	\[
	\mathcal K_D^\circ
	:=
	\{x:x_t>1\text{ for }t\in D,
	x_s<1\text{ for }s\in\mathcal N(D)\}.
	\]
	Then
	$\{u^{-1}X_{D^{\oplus 1}}\in\mathcal K_D^\circ\}
	\subset \{C_u(o)=D\}\subset
	\{u^{-1}X_{D^{\oplus 1}}\in\mathcal K_D\}$.
	The set $\mathcal K_D^\circ$ is dense in $\mathcal K_D$: one can
	approximate a boundary point by increasing its coordinates on $D$ and
	decreasing those on $\mathcal N(D)$.  Since $Q_D$ is continuous,
	\[
	\inf_{x\in\mathcal K_D^\circ}Q_D(x)
	=
	\inf_{x\in\mathcal K_D}Q_D(x)=I_D.
	\]
	The Gaussian large deviation upper and lower bounds therefore give
	\begin{equation}\label{eq:ld-one-shape-rate}
		\lim_{u\to\infty}
		-\frac1{u^2}\log\P\{C_u(o)=D\}=I_D.
	\end{equation}
	Because $\mathcal C_k^{\rm root}$ is finite and $w_k(u)$ is the sum of
	the corresponding shape probabilities, \eqref{eq:ld-one-shape-rate}
	implies
	\[
	\lim_{u\to\infty}
	-\frac1{u^2}\log w_k(u)
	=\min_{D\in\mathcal C_k^{\rm root}}I_D=I_k.
	\]
	Since $k$ is fixed, the
	second equality in \eqref{eq:ld-weight-high-u} follows from \eqref{eq:ld-weight-sandwich}.
	
	It remains to determine the normalizing rates.  Define
	\begin{equation*}
		q_u:=\P(X_o>u),
		\qquad
		b_u:=\P\!\left(
		X_o>u,\ \max_{s\in\mathcal N(o)}X_s>u
		\right).
	\end{equation*}
	If $X_o>u$ and every neighbor of $o$ is at most $u$, then $\{o\}$ is a
	singleton cluster and $o$ is both its root and its peak.  Therefore,
	\begin{equation}\label{eq:ld-normalizer-squeeze}
		q_u-b_u\leq\rho_u\leq q_u,
		\qquad
		q_u-b_u\leq\rho_u^{\rm peak}\leq q_u.
	\end{equation}
	Positive definiteness and the finiteness of $\mathcal N(o)$ imply $r_*=\max_{s\in\mathcal N(o)}
	\operatorname{Corr}(X_o,X_s)<1$.
	For each $s\in\mathcal N(o)$,
	\[
	\P(X_o>u,X_s>u)
	\leq\P(X_o+X_s>2u)
	\leq
	\exp\!\left\{-\frac{u^2}
	{\sigma^2\{1+\operatorname{Corr}(X_o,X_s)\}}\right\}.
	\]
	A union bound and the Gaussian marginal tail rate consequently yield
	\begin{equation*}
		\begin{aligned}
			\frac{b_u}{q_u}
			&\leq
			\exp\left\{
			-\frac{1-r_*}{2\sigma^2(1+r_*)}u^2+o(u^2)
			\right\}
			\to 0,\quad 
			\lim_{u\to\infty}-\frac1{u^2}\log q_u
			=\frac1{2\sigma^2}.
		\end{aligned}
	\end{equation*}
	It follows from \eqref{eq:ld-normalizer-squeeze} that
	$\rho_u\sim q_u$ and $\rho_u^{\rm peak}\sim q_u$, proving
	\eqref{eq:ld-normalizer-high-u}.  Finally, divide the weight asymptotics
	in \eqref{eq:ld-weight-high-u} by their corresponding normalizers to obtain
	\eqref{eq:ld-pmf-high-u}.
\end{proof}


\begin{corollary}[Concentration on singleton clusters]
	\label{cor:ld-singleton}
	Under the assumptions of Theorem~\ref{thm:ld-high-threshold}, we have
	\begin{equation}\label{eq:ld-singleton-rate}
		I_1=\frac1{2\sigma^2},
		\qquad
		\P(S_u=1)\to 1,
		\qquad
		\P(S_u^{\rm peak}=1)\to 1.
	\end{equation}
	For every fixed $k\geq2$,
	\begin{equation}\label{eq:ld-nonsingleton-rate}
		I_k-\frac1{2\sigma^2}
		\geq
		\frac{1-r_*}{2\sigma^2(1+r_*)}>0.
	\end{equation}
\end{corollary}

\begin{proof}
	The singleton event in \eqref{eq:ld-normalizer-squeeze} has probability
	$q_u-b_u\sim q_u$.  It follows that its rate is $1/(2\sigma^2)$ and that
	both normalized singleton probabilities tend to one.  If $k\geq2$, a
	connected size-$k$ cluster containing $o$ contains an exceedance neighbor
	of $o$.  Thus $w_k(u)\leq b_u$, and the pair bound in the proof of
	Theorem~\ref{thm:ld-high-threshold} gives
	\[
	I_k\geq\frac1{\sigma^2(1+r_*)}.
	\]
	Subtracting the normalizer rate $1/(2\sigma^2)$ proves
	\eqref{eq:ld-nonsingleton-rate}.
\end{proof}

\begin{remark}
	A singleton cluster is an isolated exceedance: it consists of one
	site \(t\) such that \(X_t>u\), while all neighbors of \(t\) are at
	or below \(u\). Thus, \(\P(S_u=1)\to1\) means that the size of a
	typical finite cluster converges in distribution to one. It does not
	mean that the entire excursion set contains no non-singleton
	clusters. The same interpretation applies to the peak-based
	distribution.
\end{remark}

\begin{example}[One-dimensional white noise in both regimes]
	\label{example:ld-white-noise}
	Let $p_u=F(u)$ and $q_u=1-p_u\in(0,1)$, as in
	Examples~\ref{example:WN_1D} and \ref{example:peak-GWN}.  The exact law
	satisfies $\P(S_u=k)=p_uq_u^{k-1}$, while for $k\geq2$ the peak-based law is
	proportional to $(k+1)q_u^k$.  Hence, for fixed $u$, both distributions have
	the exact rate
	\[
	-\lim_{k\to\infty}\frac1k\log\P(S_u=k)
	=
	-\lim_{k\to\infty}\frac1k\log\P(S_u^{\rm peak}=k)
	=-\log q_u.
	\]
	If, in addition, $F=\Phi$ and $k$ is fixed while $u\to\infty$, then
	$q_u=1-\Phi(u)$ satisfies $-u^{-2}\log q_u\to1/2$ and
	$1-p_u^3\sim3q_u$.  The same formulas therefore give
	\begin{equation}\label{eq:ld-white-noise-high-u}
		\lim_{u\to\infty}-\frac1{u^2}\log\P(S_u=k)
		=
		\lim_{u\to\infty}-\frac1{u^2}
		\log\P(S_u^{\rm peak}=k)
		=\frac{k-1}{2}.
	\end{equation}
	Thus the two regimes have different speeds: $k$ at a fixed threshold and
	$u^2$ at a fixed cluster size.
\end{example}

\begin{remark}[Scope of the two regimes]
	\label{remark:ld-scope}
	Theorem~\ref{thm:ld-general} and Corollary~\ref{cor:ld-gaussian} provide
	sufficient, generally conservative, bounds as $k\to\infty$ with $u$
	fixed.  Gaussianity alone is insufficient in that regime: long-range
	dependence or a percolating excursion regime may prevent exponential
	cluster-size decay.  In contrast,
	Theorem~\ref{thm:ld-high-threshold} gives the exact $u^2$-scale rate for
	each fixed $k$ under finite-dimensional nondegeneracy.  The two limits have
	different speeds, $k$ and $u^2$, and cannot be interchanged without
	additional uniform estimates.
\end{remark}
	
	\begin{figure}[t]
		\centering
		\begin{tabular}{@{}c@{\hspace{0.25in}}c@{}}
			\includegraphics[
			trim=10 20 20 0,clip,width=2.5in
			]{GRF_WN_CSD_1D_u0.5.png}
			&
			\includegraphics[
			trim=10 20 20 0,clip,width=2.5in
			]{GRF_WN_CSD_1D_u1.5.png}
			\\
			\multicolumn{2}{c}{
				(a) Gaussian white noise on $\mathbb{Z}$.
			}
			\\[1em]
			\includegraphics[
			trim=10 20 20 0,clip,width=2.5in
			]{GRF_sq_exp_cov_1D_u0.5.png}
			&
			\includegraphics[
			trim=10 20 20 0,clip,width=2.5in
			]{GRF_sq_exp_cov_1D_u1.5.png}
			\\
			\multicolumn{2}{c}{
				(b) Gaussian process on $\mathbb{Z}$ with covariance
				$C(t,s)=e^{-(t-s)^2}$.
			}
		\end{tabular}
		\caption{
			Exact and peak-based cluster size probability masses for
			(a) Gaussian white noise and
			(b) a Gaussian process with covariance
			$C(t,s)=e^{-(t-s)^2}$ on $\mathbb{Z}$. The theoretical values are compared with spatial Monte Carlo estimates.
			The left and right columns correspond to thresholds
			$u=0.5$ and $u=1.5$, respectively.		
		}
		\label{fig:sim-est-1D}
	\end{figure}
	
	\section{Numerical illustrations}\label{sec:simulation}
	We present numerical experiments illustrating the practical performance of the methods developed above for computing the exact and peak-based cluster size distributions and their corresponding weights, $w_k$ and $w_k^{\rm peak}$. For each example, we compare appropriate reference values with spatial Monte Carlo estimates obtained from independent realizations of the random field.
	
	As noted before, closed-form expressions for $w_k$, $w_k^{\rm peak}$, $\P(S_u=k)$, and $\P(S_u^{\rm peak}=k)$ are generally unavailable, particularly when $d\geq2$. Accordingly, the numerical results below use three types of reference values: analytic values when explicit formulas are available, numerical evaluations of the exact finite-dimensional formulas for moderate $k$, and independent large-window simulation estimates when direct evaluation is computationally infeasible.
	
	In Figures \ref{fig:sim-est-1D} and \ref{fig:sim-est-2D}, we use the first ten values of $w_k$ and $w_k^{\rm peak}$, their spatial Monte Carlo estimates $\widehat{w}_k$ and $\widehat{w}_k^{\rm peak}$ defined in \eqref{eq:MC2}--\eqref{eq:W2}, and the corresponding sums $\rho_u=\sum_{k=1}^\infty w_k$, $\rho_u^{\rm peak}=\sum_{k=1}^\infty w_k^{\rm peak}$ and  $\widehat\rho_K = \sum_{k=1}^K \widehat{w}_k$ over the computed range. Normalizing by these sums yields the cluster size probability mass functions. Specifically, the ratios
	\[
	\frac{w_k}{\rho_u}=\P(S_u=k) \quad \text{and} \quad \frac{w_k^{\rm peak}}{\rho_u^{\rm peak}}=\P(S_u^{\rm peak}=k)
	\]
	correspond to the theoretical exact and peak-based cluster size distributions respectively; and 
	\[
	\frac{\widehat{w}_k}{\widehat\rho_K} \quad \text{and} \quad \frac{\widehat{w}_k^{\rm peak}}{\rho_u^{\rm peak}}
	\]
	correspond to their spatial Monte Carlo estimates respectively. For each example, these distributions are also displayed graphically as histograms to facilitate visual comparison.
	
	
	\subsection{Gaussian processes on $\mathbb{Z}$}
	We consider Gaussian white noise and a stationary Gaussian process with
	covariance $C(t,s)=e^{-(t-s)^2}$.  For each model and threshold, we generated $M_1=10{,}000$ independent realizations on a simulation window $V$ of length $200$ and scored the centered subwindow $\Lambda$ of length $100$ with a $K=50$ guard band. 
	
	For white noise, set $p=\Phi(u)$ and $q=1-p$.  As shown in Examples \ref{example:WN_1D} and \ref{example:peak-GWN}, the reference values are
	available analytically:
	\begin{align*}
		w_k&=p^2q^k,& \rho_u&=pq,\\
		w_1^{\rm peak}&=p^2q,&
		w_k^{\rm peak}&=\frac{k+1}{3}p^2q^k\quad(k\geq2),&
		\rho_u^{\rm peak}&=\frac{1-p^3}{3}.
	\end{align*}
	For the correlated process,  theoretical values of $w_k$, $\rho_u$, $w_k^{\rm peak}$ and $\rho_u^{\rm peak}$ were computed from \eqref{eq:w_1D}, \eqref{eq:w_1D_sum}, and the numerator and denominator of \eqref{eq:peak-GP}, respectively. 
	
	Figure \ref{fig:sim-est-1D} displays histogram comparisons. In both examples the empirical distributions closely match theory. Moreover, as the threshold increases from $u=0.5$ to $u=1.5$, the peak-based cluster size distribution approach the exact cluster size distribution, confirming the high-threshold approximation discussed in Section \ref{sec:high-u}. 
	
		\begin{figure}[t]
		\centering
		\begin{tabular}{@{}c@{\hspace{0.25in}}c@{}}
			\includegraphics[
			trim=10 20 20 0,clip,width=2.5in
			]{GRF_CSD_4N_2D_u0.5.png}
			&
			\includegraphics[
			trim=10 20 20 0,clip,width=2.5in
			]{GRF_CSD_4N_2D_u1.5.png}
			\\
			\multicolumn{2}{c}{
				(a) Nearest-neighbor connectivity on $\Z^2$.
			}
			\\[1em]
			\includegraphics[
			trim=10 20 20 0,clip,width=2.5in
			]{GRF_CSD_8N_2D_u0.5.png}
			&
			\includegraphics[
			trim=10 20 20 0,clip,width=2.5in
			]{GRF_CSD_8N_2D_u1.5.png}
			\\
			\multicolumn{2}{c}{
				(b) Moore connectivity on $\Z^2$.
			}
		\end{tabular}
		\caption{ 
			Exact and peak-based cluster size probability masses for a Gaussian field on $\Z^2$ with covariance
			$C(t,s)=e^{-\|t-s\|^2}$ under
			(a) nearest-neighbor connectivity and 
			(b) Moore connectivity.
			The theoretical values are compared with spatial Monte Carlo estimates. The left and right columns correspond to thresholds
			$u=0.5$ and $u=1.5$, respectively.		
		}
		\label{fig:sim-est-2D}
	\end{figure}
	
	
	
	\subsection{Gaussian fields on $\mathbb{Z}^2$}
	We next consider the centered stationary Gaussian field on $\Z^2$ with
	\[
	\operatorname{Cov}(X_s,X_t)=e^{-\|t-s\|^2}.
	\]
	Using the same field and thresholds under nearest-neighbor and Moore
	adjacency. The four comparisons
	are collected in one display in Figure \ref{fig:sim-est-2D}.
	
	In general, theoretical values of $w_k$, $w_k^{\rm peak}$ and $\rho_u^{\rm peak}$ are computed from \eqref{eq:w_d}, \eqref{eq:w-peak} and \eqref{eq:peak-intensity-local}, respectively. Because exact numerical evaluation becomes costly for large cluster sizes, we compute theoretical values exactly for $k\le 6$ and estimate larger $k$ by empirical values. Specifically, we simulate the field on an integer lattice of size $3,000\times 3,000$. The empirical estimators of $w_k$, for $k\ge 7$, are provided in \eqref{eq:w-empirical}. Similarly, the empirical estimators for the peak-based quantities are given by \eqref{eq:w-peak-empirical}.
	
	The Monte Carlo estimates use $M_2=15{,}000$ independent realizations on a
	$200\times200$ simulation window $V$, with scores computed on the centered
	$100\times100$ window $\Lambda$ and a $K=50$ guard band.
	
	 The results in Figure \ref{fig:sim-est-2D} show that, when moving from nearest-neighbor to Moore connectivity, or from one-dimensional to two-dimensional settings, small clusters occur less frequently while larger clusters become more prevalent. This behavior is a direct consequence of the increased degree of connectivity, which promotes the formation of larger connected components. 
	
	Across all examples presented in Figures \ref{fig:sim-est-1D} and \ref{fig:sim-est-2D}, the Monte Carlo estimated exact and peak-based distributions closely match their theoretical counterparts. As the threshold increases from $u=0.5$ to $u=1.5$, the peak-based cluster size distribution converges toward the exact distribution, consistent with the high-threshold approximation in Section \ref{sec:high-u}.
	
	%%%%%%%%%%%%%%%%%%%%%%%%%%%%%%%%%%%%%%%%%%%%%%
	%% Support information, if any,             %%
	%% should be provided in the                %%
	%% Acknowledgements section.                %%
	%%%%%%%%%%%%%%%%%%%%%%%%%%%%%%%%%%%%%%%%%%%%%%
	
	
	
	\begin{funding}
		The authors were supported by NSF Grant DMS-2220523 and Simons Foundation Collaboration Grant \#854127. 
	\end{funding}
	
	%%%%%%%%%%%%%%%%%%%%%%%%%%%%%%%%%%%%%%%%%%%%%%
	%% Funding information, if any,             %%
	%% should be provided in the                %%
	%% funding section.                         %%
	%%%%%%%%%%%%%%%%%%%%%%%%%%%%%%%%%%%%%%%%%%%%%%
	
	
	
	%%%%%%%%%%%%%%%%%%%%%%%%%%%%%%%%%%%%%%%%%%%%%%
	%% Supplementary Material, including data   %%
	%% sets and code, should be provided in     %%
	%% {supplement} environment with title      %%
	%% and short description. It cannot be      %%
	%% available exclusively as external link.  %%
	%% All Supplementary Material must be       %%
	%% available to the reader on Project       %%
	%% Euclid with the published article.       %%
	%%%%%%%%%%%%%%%%%%%%%%%%%%%%%%%%%%%%%%%%%%%%%%
	
	%%%%%%%%%%%%%%%%%%%%%%%%%%%%%%%%%%%%%%%%%%%%%%%%%%%%%%%%%%%%%
	%%                  The Bibliography                       %%
	%%                                                         %%
	%%  imsart-???.bst  will be used to                        %%
	%%  create a .BBL file for submission.                     %%
	%%                                                         %%
	%%  Note that the displayed Bibliography will not          %%
	%%  necessarily be rendered by Latex exactly as specified  %%
	%%  in the online Instructions for Authors.                %%
	%%                                                         %%
	%%  MR numbers will be added by VTeX.                      %%
	%%                                                         %%
	%%  Use \cite{...} to cite references in text.             %%
	%%                                                         %%
	%%%%%%%%%%%%%%%%%%%%%%%%%%%%%%%%%%%%%%%%%%%%%%%%%%%%%%%%%%%%%
	
	%% if your bibliography is in bibtex format, uncomment commands:
	%\bibliographystyle{imsart-number} % Style BST file (imsart-number.bst or imsart-nameyear.bst)
	%\bibliography{bibliography}       % Bibliography file (usually '*.bib')
	
	%% or include bibliography directly:
	\begin{thebibliography}{4}
		%%
		\bibitem{Adler81}
		Adler, R. J. (1981). {\it The Geometry of Random Fields}. Wiley,
		New York.
		
		\bibitem{AT07}
		Adler, R. J. and Taylor, J. E. (2007). {\it Random Fields and Geometry}. Springer, New York.
		
		\bibitem{Bansal:2018}
		Bansal, R. and Peterson, B. S. (2018). Cluster-level statistical inference in fMRI datasets: The unexpected behavior of random fields in high dimensions. {\it Magn Reson Imaging},
		{\bf 49},  101--115.
		
		\bibitem{Bardeen:1985}
		Bardeen, J. M., Bond, J. R., Kaiser, N. and Szalay A. S. (1985). The statistics of peaks of Gaussian random fields. {\it Astrophys. J.},
		{\bf 304},  15--61.
		
		\bibitem{Benjamini:2019}
		Benjamini, Y., Taylor, J. and Irizarry, R. A. (2019). Selection-corrected statistical inference for region detection with high-throughput assays. \emph{Journal of the American Statistical Association}, {\bf 114}, 1351--1365.
		
		\bibitem{Cheng:2020}
		Cheng, D., Cammarota, V., Fantaye, Y., Marinucci, D. and Schwartzman, A. (2020). Multiple testing of local maxima for detection of peaks on the (celestial) sphere. {\it Bernoulli}, {\bf 26},  31--60.
		
		\bibitem{CS15}
		Cheng, D. and Schwartzman, A. (2015). Distribution of the height of local maxima of Gaussian random fields. {\it Extremes}, {\bf 18}, 213--240.
		
		\bibitem{CS18}
		Cheng, D. and Schwartzman, A. (2018). Expected number and height distribution of critical points of smooth isotropic Gaussian random fields. {\it Bernoulli}, {\bf 24}, 3422--3446.
		
		\bibitem{CS17}
		Cheng, D. and Schwartzman, A. (2017). Multiple testing of local maxima for detection of peaks in random fields. {\it Ann. Stat.}, {\bf 45}, 529--556.
		
		\bibitem{Chiles:2012}
		Chiles, J.-P. and Delfiner, P. (2012). {\it Geostatistics: Modeling Spatial Uncertainty}. Wiley, New York.
		
		\bibitem{Chumbley:2010}
		Chumbley, J. R., Worsley, K., Flandin, G. and Friston, K. J. (2010). Topological FDR for neuroimaging.
		\emph{Neuroimage}, {\bf 49}, 3057--3064.
		
		\bibitem{Cooley:2007}
		Cooley, D., Nychka, D. and Naveau, P. (2007). Bayesian spatial modeling of extreme precipitation return levels. \emph{J. Am. Statist. Assoc.}, {\bf 102}, 824--840.
		
		%\bibitem{CL67}
		%Cram\'er, H. and Leadbetter, M. R. (1967). {\it Stationary and Related Stochastic Processes:
			%Sample Function Properties and Their Applications}. Wiley, New York.
		
		\bibitem{Davison:2012}
		Davison, A. C., Padoan, S. A. and Ribatet, M. (2012). Statistical modeling of spatial extremes. \emph{Statistical Science}, {\bf 27}, 161--168.
		
		%\bibitem{Friston:1994}
		%Friston, K. J.,  Holmes, A. P.,  Worsley, K. J., Poline, J. P., Frith, %C. D. and Frackowiak, R. S. J. (1994). Statistical parametric maps in %functional imaging: A general linear approach. \textit{Human Brain %Mapping}, \textbf{2}, 189--210.
		
		\bibitem{Friston:1994}
		Friston, K. J.,  Worsley, K. J., Frackowiak, R. S., Mazziotta, J. C. and Evans, A. C.  (1994). Assessing the significance of focal activations using their spatial extent. \textit{Human Brain Mapping}, \textbf{1(3)}, 210--20.
		
		\bibitem{Fondeville:2018}
		de Fondeville, R. and Davison, A. C. (2018). High-dimensional peaks-over-threshold inference. \textit{Biometrika}, \textbf{105}, 575--592.
		
		\bibitem{Goeman:2023}
		Goeman, J. J., Gorecki, P., Monajemi, R., Chen, X., Nichols, T. E. and Weeda, W. (2023). Cluster extent inference revisited: quantification and localisation of brain activity. \textit{Journal of the Royal Statistical Society Series B}, \textbf{85}, 1128--1153.
		
		\bibitem{Huser:2014}
		Huser, R. and Davison, A. C. (2014). Space–time modelling of extreme events. \textit{JRSSB}, \textbf{76}, 439--461.
		
		\bibitem{Lindgren82}
		Lindgren, G. (1982). Wave characteristics distributions for Gaussian waves -- wave-length, amplitude and steepness. {\it Ocean Engng.}, {\bf 9},  411--432.
		
		\bibitem{MP11}
		Marinucci, D. and Peccati, G. (2011). {\it Random Fields on the Sphere. Representation, Limit Theorems and Cosmological Applications}. Cambridge University Press.
		
		\bibitem{Poline:1997}
		Poline, J. B., Worsley, K. J., Evans, A. C. and Friston, K. J. (1997). Combining spatial extent and peak intensity to test for activations in functional imaging. \emph{Neuroimage}, {\bf 5}, 83--96.
		
		\bibitem{Vanmercke:2010}
		Vanmercke, E. (2010). \textit{Random Fields: Analysis and Synthesis.} World Scientific.
		
		\bibitem{Tan:2021}
		Tan, X., Wu, X., and Liu, B. (2021). Global changes in the spatial extents of precipitation extremes. \emph{Environmental Research Letters}, {\bf 16}, 054017.
		
		\bibitem{Worsley:1996a}
		Worsley, K. J., Marrett, S., Neelin, P., Vandal, A. C., Friston, K. J. and Evans, A. C. (1996). A unified statistical approach for determining significant signals in images of cerebral activation. \emph{Human Brain Mapping}, {\bf 4}, 58--73.
		
		\bibitem{Taylor:2007}
		Taylor, J. E. and Worsley, K. J. (2007). Detecting sparse signals in random fields, with an application to brain mapping. \emph{Journal of the American Statistical Association}, {\bf 102}, 913--928.
		
		\bibitem{Zhang:2009}
		Zhang, H.,  Nichols, T. E. and Johnson, T. D. (2009). Cluster mass inference via random field theory. \emph{Neuroimage}, {\bf 44}, 51--61.
		
		\bibitem{Zhong:2025}
		Zhong, P., Brunner, M., Opitz, T. and Huser, R. (2025). Spatial modeling and future projection of extreme precipitation extents. \emph{Journal of the American Statistical Association}, {\bf 120}, 80--95.
	\end{thebibliography}
	
\end{document}
